% !Mode:: "TeX:UTF-8"

\section{课题来源及研究的目的和意义}

\subsection{课题来源与研究背景}
本课题来源于课题组轮式双臂机器人在PCBA检测工位上下料中的真机部署需求。前期工作已完成VR示范采集、OpenPI \(\pi_{0.5}\)微调与机器人执行接口联调；在此基础上，工位停靠残差对精细操作可靠性的影响及其抑制方法构成本课题拟解决的工程问题。

工业机器人是面向工业环境设计、可自动控制、可重复编程的多关节机械装置，凭借高精度、高节拍和长时间稳定运行的优势，长期承担焊接、喷涂、搬运、装配等重复性作业，是制造业自动化的核心装备。国际机器人联合会（International Federation of Robotics, IFR）发布的2025年世界机器人报告显示，2024年全球工业机器人新装机量达54.2万台，较十年前增长一倍以上，全球在役存量达466万台；其中中国新装机29.5万台，占全球的54\%，在役存量突破200万台，连续多年保持全球最大的工业机器人市场\cite{ifr2025worldrobotics}。电子电气行业是工业机器人的主要应用行业之一，电路板组件的上下料、插件与检测等环节对机器人替代人工存在持续需求。然而，传统工业机器人依赖示教编程或离线编程，在结构化环境中沿预设轨迹重复运动，其可靠性建立在工件位姿、夹具和工艺流程高度固定的前提之上\cite{yuan2025embodiedmanufacturing,zhang2026eiir}。当产线面临多品种、小批量、快速换型的需求，或工件位置、来料状态存在不确定性时，机器人需要重新示教、重新标定甚至重新设计工装，部署周期与成本显著上升；固定安装的机器人还需要依靠传送带或专用输送设备衔接不同工位，难以独立承担跨工位的物料流转任务。

与传统工业机器人依赖精确编程不同，具身智能（Embodied Intelligence）强调智能体通过物理载体与环境的动态交互，在感知、决策与行动的闭环中持续学习和进化，从而突破静态数据训练模型的局限，获得更强的环境适应性与泛化能力\cite{zheng2026embodiedtrend}。近年来，在互联网规模数据上预训练的大语言模型与视觉--语言模型（Vision-Language Model, VLM）展现出开放词汇视觉识别、常识推理与任务规划能力，推动机器人基础模型成为具身智能的技术核心\cite{firoozi2025foundation}。在操作技能层面，RT-1与RT-2将Transformer架构和网络规模的视觉语言知识引入真机控制，使机器人能够从大规模示范数据中学习语言条件下的操作技能\cite{brohan2022rt1,brohan2023rt2}；Open X-Embodiment汇集跨本体、跨机构的机器人数据，验证了跨机器人数据的正向迁移\cite{openx2024rtx}；视觉--语言--动作（Vision-Language-Action, VLA）模型以预训练VLM为主干、以动作专家生成连续动作，在开放环境中对新物体、新场景与新指令表现出泛化能力\cite{kim2024openvla,black2024pi0,black2025pi05}。与此同时，低成本遥操作与动作分块模仿学习使双臂精细操作可以从数十至数百条人类示范中直接学得\cite{zhao2023act,fu2024mobilealoha}；移动操作机器人因兼具移动、规划与操作能力，成为具身智能优先选择的硬件载体\cite{lan2024mobilemanipulator}。具身智能机器人由此呈现出与传统工业机器人不同的技术形态：通过数据驱动策略降低对逐点示教和完全结构化环境的依赖，并在闭环感知中生成动作。

将具身智能引入工业制造，使机器人在产线上具备自主感知、决策与操作能力，正在成为工业机器人技术演进的新方向。2025年《政府工作报告》首次将具身智能列为国家重点培育的未来产业\cite{gov2025report}；智能制造也正经历从基于规则的自动化制造、数据驱动的数字化智能制造向具身智能赋能的智能制造的演进，具身智能在生产制造、仓储物流、检测维护和人机协作等环节具有直接的赋能作用\cite{yuan2025embodiedmanufacturing,ren2025embodiedsmartmanufacturing}。然而，工业场景要求机器人更高效、更准确、更可靠、更安全地完成作业，对具身智能提出了不同于开放家庭场景的挑战\cite{zhang2026eiir}：操作对象往往是对准公差小、失效代价高的精密工件，产线对节拍与合格率有明确约束，策略还须在多台机器人之间稳定部署。本课题前期调试反复遇到由移动底盘停靠偏差引起的精细操作失效，而现有名义停靠示范不覆盖这类状态。如何在固定PCBA型号、工装几何与环境条件下，量化并提高策略对工位停靠残差的容忍能力，是本课题的研究出发点；来料位姿、光照、工装和产线差异不作为本次实验变量。

本课题面向的印制电路板组件（Printed Circuit Board Assembly, PCBA）检测工位上下料正是这样一个典型场景，它是电子制造产线中重复性高、节拍受限的环节。一次完整的上下料由四个物理工位上的五个操作子任务组成：从上料百叶车夹取待测PCBA、移动至扫码工位完成扫码登记、从检测台取出已完成检测的上一块PCBA、将当前PCBA放入检测台、最后将上一块PCBA放入下料百叶车；检测台的取出与放入共用一次停靠。工位间移动由导航系统完成，不计入VLA动作阶段；进出百叶车槽位、检测台取放等接触操作属于精细阶段，机械臂在工位内的接近、搬运与撤离属于自由阶段。本课题所用轮式双臂机器人由移动底盘、四自由度腰部、两条七自由度机械臂、二自由度头部与末端夹爪组成，配置头部相机和左右腕部相机，并由激光雷达定位与传统导航算法在工位间移动。前期观测表明实际停靠同时存在厘米量级的位置残差和度量级的航向残差，其统计分布及对操作成功率的影响将在正式实验中按统一协议测量，而不由调试现象代替统计结论。

面向这类多阶段双臂任务，以遥操作示范驱动的模仿学习已成为主流策略学习路线。动作分块Transformer（Action Chunking with Transformers, ACT）通过一次预测一段动作序列，提出了降低长时序行为克隆有效决策时域的动作分块思想\cite{zhao2023act}；流匹配VLA模型进一步继承连续动作块表示，并借助预训练视觉语言表征支持多任务与语言条件控制\cite{black2024pi0,black2025pi05}。本课题已在上述平台上完成“虚拟现实（Virtual Reality, VR）遥操作数据采集--基于OpenPI的\(\pi_{0.5}\)微调--真机执行”流水线，能够在名义停靠位姿下完成一次从观测到动作执行再到新观测的接口闭环；该结果只说明链路可运行，不作为稳定成功率结论。ACT为动作分块策略提供了方法基础，本文在\(\pi_{0.5}\)上研究工位停靠残差。前期调试已经观察到停靠偏差条件下的进槽、落台失败及机械挤压风险，但失败概率随残差的变化仍待分档重复实验量化。

本研究将该现象归纳为一条误差传递链，如图\ref{fig:embodied-control-loop}所示。首先，导航停靠产生平面位置与航向残差；其次，位置残差近似等幅改变目标相对位置，航向残差则经目标到基座的距离转化为横向误差，两者在进槽、落台等精细阶段与允许间隙直接竞争；最后，系统虽在动作块重规划边界形成视觉闭环，但一次提交的\(s\)个控制步在策略层面开环执行，其间的新观测不进入策略，错误动作因而可能持续至下一次重规划。该链条上的三个环节对应三种不同性质的缺陷：状态输入缺少工位相对位姿信息、策略级开环区间与阶段精度需求不匹配、运行时缺少失败检测与恢复机制。

\begin{figure}[htbp]
  \centering
  \begin{tikzpicture}[
  font=\small,
  >={Stealth[length=2.2mm]},
  flow/.style={draw=blue!60!black,rounded corners=2pt,fill=blue!7,
    minimum height=13mm,text width=27mm,align=center,inner sep=3pt},
  plant/.style={draw=orange!75!black,rounded corners=2pt,fill=orange!10,
    minimum height=13mm,text width=31mm,align=center,inner sep=3pt},
  constraint/.style={draw=gray!65,rounded corners=5pt,fill=gray!8,
    text width=30mm,align=center,inner sep=2.5pt},
  mainarrow/.style={->,thick,blue!70!black},
  feedback/.style={->,thick,orange!80!black,rounded corners=3pt},
  influence/.style={->,densely dashed,gray!70}
]
  \node[flow] (sense) {\textbf{多模态感知}\\视觉、力/触觉\\与本体状态};
  \node[flow,right=6mm of sense] (policy) {\textbf{具身策略}\\状态表示\\与动作决策};
  \node[flow,right=6mm of policy] (safety) {\textbf{安全/柔顺接口}\\限幅、滤波\\与低层控制};
  \node[plant,right=6mm of safety] (plant) {\textbf{机器人--工件--环境}\\装配、插接\\与约束运动};

  \draw[mainarrow] (sense) -- node[above,font=\footnotesize]{观测} (policy);
  \draw[mainarrow] (policy) -- node[above,font=\footnotesize]{动作} (safety);
  \draw[mainarrow] (safety) -- node[above,font=\footnotesize]{指令} (plant);
  \draw[feedback] (plant.south) -- ++(0,-8mm) -|
    node[pos=0.55,below,font=\footnotesize]{接触结果与环境变化反馈} (sense.south);

  \node[constraint,above=9mm of sense] (precision) {\textbf{精度约束}\\位姿误差·尺寸公差};
  \node[constraint,above=9mm of policy] (cycle) {\textbf{节拍约束}\\控制频率·推理延迟};
  \node[constraint,above=9mm of safety] (risk) {\textbf{安全约束}\\接触力·碰撞·失效代价};
  \node[draw=gray!50,densely dashed,rounded corners=3pt,fit=(precision)(cycle)(risk)(plant),
    inner sep=4mm,label={[font=\footnotesize,text=gray!70!black]above right:工业部署边界}] {};
  \draw[influence] (precision) -- (sense);
  \draw[influence] (cycle) -- (policy);
  \draw[influence] (risk) -- (safety);
\end{tikzpicture}

  \caption{PCBA上下料任务中的模仿学习策略闭环与工位停靠残差传递链}
  \label{fig:embodied-control-loop}
\end{figure}

针对上述三个环节，本课题提出基于工位停靠残差传递链的分层鲁棒增强框架。该框架由停靠残差条件化模仿学习方法、阶段自适应动作块控制方法和预测一致性驱动的运行时容错方法构成：条件化学习方法将定位系统估计的停靠残差作为\(\pi_{0.5}\)的状态条件，并以随机停靠示范扩大条件空间的数据覆盖；自适应控制方法按任务阶段调节动作块执行时域，在精细阶段缩短开环区间，并以实时动作分块（Real-Time Chunking, RTC）连续衔接前后动作块；运行时容错方法利用RTC重叠区中新旧动作块的归一化一致性残差检测失败风险，并以固定模板执行恢复与重试。三项方法均不改变策略网络主体结构，分别作用于数据与状态、推理执行和运行时监控，可依次叠加并在同一残差注入协议下比较。这一从已贯通\(\pi_{0.5}\)数据--推理--动作执行接口链路的真机需求出发、沿残差传递链构建分层增强机制的研究情境，构成本课题的直接来源。

\subsection{研究目的与意义}
本研究的总体目的是，面向PCBA上下料这一跨多工位、含精细对准阶段的轮式双臂任务，在已贯通的OpenPI \(\pi_{0.5}\)数据--推理--动作执行接口链路基础上，沿工位停靠残差传递链建立可逐层叠加、可统一验证的策略鲁棒性改进方法。\(\pi_{0.5}\)是本课题唯一的工程部署与实验主线，具体目的包括以下三个方面。

\begin{enumerate}[label=\textup{(\arabic*)},leftmargin=2.5em,itemsep=0.25em,topsep=0.25em]
  \item 停靠残差条件化模仿学习：将激光雷达定位估计\(\widehat{\boldsymbol\delta}^{\mathrm{loc}}=(\widehat\delta_x^{\mathrm{loc}},\widehat\delta_y^{\mathrm{loc}},\widehat\delta_\theta^{\mathrm{loc}})\)作为显式状态输入，并在示范采集时用导航指令\(\boldsymbol\delta^{\mathrm{cmd}}\)对\(x\)、\(y\)各施加\(\pm3\)厘米、对航向施加\(\pm3^\circ\)的随机停靠，使策略学习“估计残差--动作修正”的对应关系；通过固定停靠、随机停靠和随机停靠加残差输入三组消融，分离数据覆盖与显式条件两种作用。
  \item 阶段自适应动作块控制：在不重新训练模型的前提下，按自由空间与精细操作两类阶段设置不同的执行时域，精细阶段仅执行动作块前部即重新推理，并以RTC冻结已确定动作、引导重叠区和生成后续动作，以缩短开环执行区间并保持跨块连续性；与固定执行时域方式在成功率和推理时延两方面比较。
  \item 预测一致性驱动的运行时容错：以RTC重叠区中新旧动作块的归一化残差定义跨块不一致度，以独立成功标定集的全回合统计量确定阈值，实现无需失败训练数据和额外模型训练的失败检测，并配合回退至安全高度、整链最多重试一次的模板化恢复；先用冻结C2-A轨迹在C3-D离线影子模式下评价回合召回率、误报率和真阳性提前量，再在阈值冻结后以独立C3-R回合评价条件恢复成功率、最终整链成功率及分项时延。
\end{enumerate}

图\ref{fig:research-entry-route}给出了本课题的研究切入路线。已贯通并完成单回合功能验证的OpenPI \(\pi_{0.5}\)接口链路提供了功能基线与真机测试条件；真机部署中的失败观察引出工位停靠残差传递链；分层鲁棒增强框架中的三个方法分别对应链条的三个环节，并在统一的残差注入协议、五子任务链和四台机器人条件下按“C0基线方法--C1条件化学习--C2叠加自适应控制--C3完整框架”的顺序验证。验证结果还将反向修正示范采集规范与部署参数。

\begin{figure}[htbp]
  \centering
  \begin{tikzpicture}[
  font=\small,
  >={Stealth[length=2.2mm]},
  stage/.style={draw=blue!60!black,rounded corners=2pt,fill=blue!7,
    minimum height=21mm,text width=24mm,align=center,inner sep=2.5pt},
  validate/.style={draw=orange!75!black,rounded corners=2pt,fill=orange!10,
    minimum height=21mm,text width=27mm,align=center,inner sep=2.5pt},
  issue/.style={draw=red!55!black,rounded corners=4pt,fill=red!4,
    text width=24mm,align=center,inner sep=2pt,font=\footnotesize},
  metrics/.style={draw=red!55!black,rounded corners=4pt,fill=red!4,
    text width=27mm,align=center,inner sep=2pt,font=\footnotesize},
  arrow/.style={->,thick,blue!70!black},
  feedback/.style={->,thick,densely dashed,orange!80!black,rounded corners=3pt},
  issuearrow/.style={->,densely dashed,red!65!black}
]
  \node[stage] (task) {\textbf{任务与评价协议}\\观测·动作\\约束·扰动};
  \node[stage,right=5mm of task] (prior) {\textbf{示教先验学习}\\数据采集与对齐\\$\downarrow$\\模仿策略};
  \node[stage,right=5mm of prior] (ppo) {\textbf{PPO策略优化}\\标准基线\\与候选改进};
  \node[validate,right=5mm of ppo] (validation) {\textbf{分级验证}\\仿真对比\\$\downarrow$\\安全准入→真机};

  \draw[arrow] (task) -- (prior);
  \draw[arrow] (prior) -- (ppo);
  \draw[arrow] (ppo) -- (validation);

  \node[issue,below=8mm of prior] (dataissue) {数据覆盖与示教噪声};
  \node[issue,below=8mm of ppo] (learnissue) {样本效率·先验遗忘\\稀疏奖励·动作平滑};
  \node[metrics,below=8mm of validation] (metric) {成功·节拍·接触力\\成本·扰动退化·安全};

  \draw[issuearrow] (dataissue) -- (prior);
  \draw[issuearrow] (learnissue) -- (ppo);
  \draw[issuearrow] (metric) -- (validation);
  \draw[feedback] (metric.south) -- ++(0,-5mm) -|
    node[pos=0.55,below,font=\footnotesize]{多维评价反馈} (task.south);
\end{tikzpicture}

  \caption{从已贯通接口链路到分层鲁棒增强框架的研究切入路线}
  \label{fig:research-entry-route}
\end{figure}
\FloatBarrier

本研究的工程意义在于形成一套能够按残差、成功率、时延和安全事件验收的部署方案。显式残差输入与随机停靠数据增强旨在相对于C0正式测得的基线扩大容忍边界；分阶段执行旨在缩短精细阶段的策略级开环区间；失败检测与模板恢复旨在报警后限制错误动作持续时间。三项方法是否减少失败和人工接管由统一真机试验判定，其结论限定于相同PCBA型号和工装条件下的四台机器人，并为同类工位迁移提供接口与实验依据。

本研究的方法意义在于，将“模仿学习策略对位姿扰动不鲁棒”这一笼统判断分解为误差传递链上三个可独立度量、可分层验证的问题。三项方法均围绕OpenPI \(\pi_{0.5}\)实现：残差输入只扩展状态维度，分阶段执行与RTC只调整推理侧逻辑，失败检测复用动作块滚动推理的重叠区，无需失败训练数据和额外检测模型。上述属性可以由实现改动、训练数据和运行时开销直接核验；其对其他流匹配策略和移动操作任务的可迁移性不在本课题中作性能承诺。

综上，本课题的研究必要性源于模仿学习策略从实验室名义条件走向产线部署时，对工位停靠残差、执行粒度和运行时失败三方面缺乏系统处理的现状。为明确相关问题的已有解决路径与适用边界，后续按研究对象与方法类型对国内外相关研究进行梳理与评述。

\section{国内外研究现状及分析}

\subsection{工业操作策略学习研究现状}

本研究关注的对象是工业场景中由机器人闭环完成的抓取、放置与插接类操作，其特点是对准公差小、失效代价高、节拍要求明确。早期学习方法主要面向接触丰富的装配任务，通过强化学习或力控制在交互中学习插入与对准动作。Levine等将引导策略搜索用于接触操作，以局部轨迹优化为神经网络策略提供监督信号\cite{levine2015gps}；Inoue等在高精度装配中采用深度强化学习，使策略根据状态反馈学习探索与插入动作\cite{inoue2017assembly}；残差强化学习保留既有控制器，仅学习基础控制之外的修正量\cite{johannink2019residual}。对接触丰富操作的综述指出，样本效率、奖励建模、仿真到真机迁移与安全约束是该类方法进入产线的主要障碍\cite{elguea2023review}。上述路线依赖真机交互或高保真仿真，在多阶段、跨工位的产线任务中难以直接铺开。

基于示范的模仿学习通过专家轨迹建立观测到动作的映射，能够在不进行大范围随机探索的条件下为机器人提供可执行的策略。早期方法通常将任务分解为接近、搜索、对准和插入等阶段，再以动态或概率运动原语显式表示轨迹分布\cite{yang2019adaptive,qin2019precision,zang2023promp}。这类方法对少量数据较为友好，且中间变量便于解释，但依赖人工分段与任务坐标系先验，难以复用到对象、场景和指令持续变化的任务中。RoboMimic基准对不同质量的人类示范和多类离线学习算法进行了系统比较，表明观测空间、历史信息、示范质量与模型选择都会影响策略性能\cite{mandlekar2022robomimic}，端到端行为克隆由此成为主流。

动作分块是当前端到端模仿学习的核心技术，其与低成本遥操作的结合使多阶段双臂操作可以直接从示范中学得，而无需人工设计接触模式或奖励。ACT用条件变分自编码器概括示范动作的多模态分布，并以Transformer在时刻\(t\)一次预测长度为\(H\)的动作块\cite{zhao2023act}：
\begin{equation}
  \mathbf A_t
  =
  \left[\mathbf a_t,\ldots,\mathbf a_{t+H-1}\right].
  \label{eq:action-chunk-background}
\end{equation}
动作块联合预测把逐步决策转化为较短的块级决策，从而缩短长时序行为克隆的有效决策时域。ACT执行时对多次推理中指向同一目标时刻的重叠预测作指数加权的时序集成，以降低动作抖动；这说明执行时域决定新观测进入控制回路的频率，而重叠动作块包含可用于判断跨块连续性的信息。ACT本身是面向任务的视觉--本体状态策略，不含语言条件和预训练通用表征，因此在本文中只作为动作分块与重叠预测思想的算法背景，不作为研究模型或实验对象。

扩散策略将动作序列建模为条件去噪过程，以概率生成方式表达同一观测下的多种合理动作，并以滚动时域方式执行预测序列的近期部分后重新规划\cite{chi2023diffusion}，其总体结构如图\ref{fig:diffusion-policy-overview}所示。3D扩散策略与RISE分别以点云替代图像作为输入，提高了策略对视角与外观变化的鲁棒性\cite{ze2024dp3,wang2024rise}；航点式模仿学习则自动从示范中抽取少量关键路点，以缩短行为克隆的决策时域\cite{shi2023awe}。这些工作均把预测时域和每次执行的时域作为固定超参数：预测时域增大有利于序列一致性，但若执行时域随之增长会延长开环区间；缩短执行时域能够提高闭环反馈频率，却增加推理负担与跨块衔接需求。

\begin{figure}[htbp]
  \centering
  \includegraphics[width=\textwidth]{diffusion-policy-overview.pdf}
  \caption{扩散视觉动作策略的观测条件化与动作序列生成框架\cite{chi2023diffusion}}
  \label{fig:diffusion-policy-overview}
\end{figure}

视觉--语言--动作模型在上述视觉动作策略基础上引入语言条件和视觉语言预训练知识，使下游任务微调所需示范量显著减少，并能够在一个模型中处理多任务与语言指令\cite{black2025pi05}。RT-1以Transformer在大规模真机数据上学习语言条件策略\cite{brohan2022rt1}，RT-2把视觉语言模型直接微调为动作输出模型\cite{brohan2023rt2}；Open X-Embodiment聚合多机构、多本体数据训练RT-X模型\cite{openx2024rtx}，Octo与OpenVLA在其上分别以扩散动作头和离散动作词元建立开源通用策略\cite{octo2024octo,kim2024openvla}；DROID提供了跨场景、跨操作者的大规模真机示范数据\cite{khazatsky2024droid}。围绕双臂操作与部署效率，RDT-1B采用扩散Transformer与统一动作空间学习双臂表示\cite{liu2025rdt}，TinyVLA、RoboMamba与UniVLA分别从轻量骨干、高效序列建模和任务中心潜在动作方面降低训练与推理成本\cite{wen2025tinyvla,liu2024robomamba,bu2025univla}。

在动作表示方面，Physical Intelligence提出的\(\pi_0\)将预训练视觉语言模型与独立的动作专家结合，动作专家以流匹配目标生成连续动作块，从而在保留语言与视觉语义表征的同时适配高频、连续的机器人控制\cite{black2024pi0}。其训练过程区分覆盖多任务、多本体数据的预训练与面向下游技能的后训练：前者用于获得可迁移的通用能力，后者使用质量更高的任务示教形成稳定执行策略。ACT所代表的任务专用动作分块行为克隆说明了序列化动作输出的价值；\(\pi_0\)则把连续动作块作为VLA动作专家的生成对象，并把任务示教微调纳入通用模型的训练链路。二者在动作块表示上存在方法联系，但模型条件、训练数据与生成机制并不相同。

随后，\(\pi_0\)-FAST采用频域动作序列标记化（Frequency-space Action Sequence Tokenization, FAST），将高频连续轨迹压缩为离散动作标记，以自回归方式提高大规模VLA训练效率\cite{pertsch2025fast}。该方法与\(\pi_0\)的流匹配路线形成连续动作生成和离散动作建模两种实现。\(\pi_{0.5}\)在此基础上进一步整合两条路线：预训练阶段采用FAST离散动作标记，以便将机器人动作数据与图文、目标检测和高层语义预测等异构任务共同训练；后训练阶段加入与\(\pi_0\)一致的流匹配动作专家，以连续动作块支持实时控制\cite{black2025pi05}。在推理阶段，\(\pi_{0.5}\)先根据视觉观测与总体指令预测语言形式的高层子任务，再以该子任务为条件生成低层连续动作，从而把任务分解与动作控制纳入同一模型。由此，\(\pi\)系列形成了“\(\pi_0\)建立流匹配动作专家--\(\pi_0\)-FAST提高动作序列训练效率--\(\pi_{0.5}\)通过异构协同训练和高低层统一推理扩展开放环境泛化”的演进路径。

尽管示范驱动的策略学习大幅降低了工程门槛，工业抓放对精度的要求并未降低。IndustReal在仿真中训练取放与插装策略并迁移到真实工业机械臂，为此专门设计了几何奖励、采样课程与部署端动作积分器，以消除接触仿真误差与稳态对准误差\cite{tang2023industreal}；Sóti等针对精密抓放改进Transporter网络的平移与旋转精度，并以由粗到细的迭代推理在真实插块任务上验证，说明毫米级的放置偏差即决定任务成败\cite{soti2023precisepick}。因此，模仿学习策略进入工业抓放场景时，需要在示范驱动的便利性与精细阶段的对准精度之间建立明确的处理机制。本课题的PCBA进槽与落台正属于此类精细阶段。

综上，VLA模仿学习已形成基于通用模型、通过任务示范进行下游微调的技术路线，且以动作块作为重要输出形式。ACT为动作分块策略提供了方法基础，本文在已贯通接口链路的\(\pi_{0.5}\)上开展面向工位停靠残差的鲁棒性改进，重点检验数据与状态输入、动作块执行粒度和运行时监控三个环节能否与上述精度需求相匹配。

\subsection{移动操作停靠残差鲁棒性研究现状}

轮式双臂机器人在工位间移动后再执行操作，属于移动操作（Mobile Manipulation）问题，其核心困难之一是底盘定位误差向操作阶段的传递。Mobile ALOHA将ALOHA双臂系统安装在移动底盘上，以全身遥操作同步采集底盘与双臂示范，并把底盘速度纳入模仿学习的动作空间；该工作还表明与静态ALOHA数据协同训练能够显著提高移动任务的数据效率，但其移动与操作由同一策略端到端输出，未显式处理停靠残差\cite{fu2024mobilealoha}。ReLMoGen让学习策略只预测子目标，由经典运动生成器负责到达，体现了学习决策与经典运动执行分层的思路\cite{xia2021relmogen}；N\(^2\)M\(^2\)在给定末端轨迹的条件下学习底盘与躯干运动，使移动操作机器人在未知环境中保持末端运动学可达\cite{honerkamp2023n2m2}。上述工作把底盘运动纳入学习范围，但均以底盘持续运动为前提，与本课题“停靠后锁定底盘、仅由双臂操作”的产线流程不同。

对停靠后操作误差的补偿主要出现在力觉与技能组合研究中。MOMA-Force在真实移动操作接触任务中同时模仿动作与目标力/力矩，由导纳式全身控制器跟踪，直接处理底盘停靠偏差带来的末端接触误差\cite{yang2023momaforce}；WildLMa以全身控制器为底层、少量示范技能为中层、大语言模型组合技能为高层，在真实环境中完成长时序移动操作\cite{qiu2025wildlma}。Mandlekar等提出的人在回路任务与运动规划框架把自由空间搬运交给自动规划，仅对接触段采集遥操作示范，从数据采集角度区分了自由空间与精细阶段\cite{mandlekar2023hitltamp}。Ishida等针对移动操作机器人换工位后自遮挡与背景变化导致的失败，提出关注任务相关视角与图像区域的鲁棒模仿学习方法，说明停靠位姿差异会同时改变本体状态与图像观测\cite{ishida2024viewpoints}。

综合来看，在本文检索并纳入比较的代表性研究中，底盘停靠偏差主要通过力控补偿、技能组合或视觉表征处理，尚未发现同时把定位系统估计的工位停靠残差作为策略状态输入、并在示范采集中系统覆盖随机停靠的方案。从策略学习的角度看，工位停靠残差本质上是一类位姿扰动，其度量与应对可借鉴更一般的位姿扰动鲁棒性研究。

模仿学习策略对训练分布之外的观测缺乏保证，位姿扰动是工业部署中常见的分布偏移来源之一。Pumacay等提出的THE COLOSSEUM基准把对象位姿、相机位姿、光照、背景、干扰物与物理参数等扰动轴分离开来，逐轴报告策略性能随扰动强度的变化，从而把“泛化能力”转化为可分级度量的退化曲线\cite{pumacay2024colosseum}，其任务与扰动组织方式如图\ref{fig:colosseum-perturbations}所示。Jin等的案例分析进一步表明，视觉策略可能同时依赖任务相关物理线索与背景视觉属性\cite{jin2024cloudgripper}；前述Ishida等的工作则说明，移动操作中的视角与背景变化同样源于停靠位姿差异\cite{ishida2024viewpoints}。工位停靠残差改变目标在三路相机中的相对视角和基座--工位几何关系，而原14维本体状态不含后一关系，因此形成视觉分布变化与状态条件缺失的耦合问题。

\begin{figure}[htbp]
  \centering
  \includegraphics[width=0.76\textwidth]{colosseum-perturbations.pdf}
  \caption{Pumacay等提出的THE COLOSSEUM基准：任务层级与扰动维度\cite{pumacay2024colosseum}}
  \label{fig:colosseum-perturbations}
\end{figure}

提高扰动鲁棒性的主要途径是扩大训练数据对扰动的覆盖，或向策略提供扰动的显式条件。动力学随机化通过扩大仿真参数分布提高策略对模型偏差的容忍度\cite{peng2018dynamics}，其思想同样适用于真机示范采集中的初始位姿随机化；航点式模仿学习把动作锚定在少量关键路点上，减少了对精确轨迹复现的依赖\cite{shi2023awe}。在评价协议方面，CALVIN以语言条件的长时序任务链考察策略连续完成多个子任务的成功率，显式呈现误差随任务链长度的累积\cite{mees2022calvin}；LIBERO以程序化生成的任务套件评估策略的知识迁移与遗忘\cite{liu2023libero}；DROID与Open X-Embodiment提供了跨机器人、跨场景的大规模数据组织范式\cite{khazatsky2024droid,openx2024rtx}；RoboMimic则强调示范质量与训练终止规则对结论的影响\cite{mandlekar2022robomimic}。据此，本课题将训练、超参数验证、阈值标定与最终测试按采集日期和回合分组隔离，避免同一轨迹片段跨集合出现。

对本课题而言，上述工作提供了两方面依据。其一，工位停靠残差可以作为一条明确的扰动轴，以程序化注入的方式分级施加，并绘制成功率随残差幅值变化的容忍曲线；其二，“随机停靠数据增强”对应数据覆盖，“残差状态输入”对应显式条件，二者在文献中通常独立出现，其相对贡献需要通过固定停靠、随机停靠、随机停靠加残差输入三组消融来识别。这正是停靠残差条件化模仿学习方法的出发点与实验设计来源。

\subsection{动作块策略闭环执行研究现状}

动作块策略部署后，每次执行的时域决定了观测反馈的闭环频率，运行时失败检测则在任务层面闭合回路，二者共同影响真实系统的响应速度与错误动作持续时间。在执行方式上，本文比较的代表性方法体现了策略级开环区间长度与推理频率之间的权衡。ACT的消融表明块长增大总体上提高成功率，而时序集成通过对重叠预测加权平均进一步改善动作平滑性；该结果在方法层面说明重叠动作块具有连续性信息\cite{zhao2023act}。扩散策略同样采用固定的预测时域与执行时域，并指出执行时域过长会削弱对新观测的响应\cite{chi2023diffusion}。但对扩散或流匹配策略直接平均不同生成样本可能产生不可执行的中间动作，不能把ACT时序集成直接视为\(\pi_{0.5}\)的原生机制。RTC针对动作分块流策略，将推理延迟内已确定执行的动作冻结为前缀，并在重叠区引导生成过程、补全后续动作，从而在不重新训练策略的条件下实现异步连续执行\cite{black2025rtc}。OpenVLA-OFT系统比较了VLA微调中的动作解码方式、动作表示与学习目标，并把动作吞吐率作为与成功率并列的评价指标\cite{kim2025openvlaoft}；FAST通过频域动作词元压缩缩短自回归解码长度\cite{pertsch2025fast}，\(\pi_{0.5}\)则以流匹配动作专家非自回归地生成整段动作块\cite{black2025pi05}，两者都把降低单次推理成本作为连续控制的前提。

本文所比较的动作分块与滚动时域方法通常在部署时采用单一执行时域。对于PCBA上下料这类自由空间移动与精细对准交替出现的任务，长执行时域在进槽、落台等接触阶段可能过度开环，而全程采用短执行时域又会增加推理次数。自由空间阶段可以容忍较长的策略级开环执行，精细阶段则需要更短的执行区间以便及时利用新观测。在推理侧按阶段切换执行时域、以RTC处理跨块衔接，并同时报告条件成功率、推理频率和时延，尚需系统实现与评价，这是阶段自适应动作块控制方法所要验证的内容。

部分由停靠残差引起的失败可能在动作预测差异中提前显现，运行时失败检测的目标是在造成损坏之前识别这一先兆；该假设需要由带失败时刻标注的独立测试回合验证。FAIL-Detect提出只用成功示范数据进行模仿策略在线失败检测的模块化方法：从观测、视觉特征与策略动作中提取失败相关标量，以流密度估计建模不确定性，并用保形预测校准时变阈值，将失败识别表述为序列分布外检测，从而无需预先枚举失败类型或采集失败训练数据\cite{xu2025faildetect}。Liu等在交互式模仿学习部署中于共享潜空间训练动力学模型与失败分类器，通过对未来轨迹采样预测异常与失败，并结合分布外判据触发人工介入\cite{liu2024runtimemonitor}。在恢复方面，RACER以视觉语言模型监督器在线描述错误与空间修正，由语言条件策略执行恢复动作，训练数据包含自动生成的扰动与恢复轨迹\cite{dai2025racer}；REFLECT把多模态历史压缩为可解释的失败报告，并由语言规划器生成纠正方案\cite{liu2023reflect}。

多种代表性检测方案依赖额外的密度估计、分类器或恢复模块，且其信号来源与策略本身的动作生成过程相互独立。动作块滚动推理天然提供了另一类信号：前一动作块的未执行尾部与当前观测下的新动作提议在时间上存在重叠，当状态处于示范分布内时，二者应具有较高一致性；当工位停靠残差把状态推离示范分布或末端已经卡滞时，跨块残差可能增大。本文利用RTC构造流匹配策略的连续执行边界，同时在RTC引导前比较时间对齐的新旧动作块，以归一化残差定义不一致度，并用成功回合分位数标定阈值。这样既不直接平均不同流匹配样本，也无需引入额外模型或失败训练数据；该信号能否稳定先于失败上升，仍需在本课题任务中实证检验，由此形成预测一致性驱动的运行时容错方法。

\subsection{国内外研究现状简析}

上述代表性研究呈现出三点共同趋势。第一，多阶段双臂操作可以通过低成本遥操作示范与动作分块模仿学习获得，VLA预训练进一步降低了下游任务的示范需求。第二，位姿扰动是模仿学习策略部署中的重要分布偏移来源，扩大数据覆盖与提供显式条件是两条互补的鲁棒性途径，评价应报告性能随扰动强度的退化曲线而非单一成功率。第三，动作块执行方式与运行时监控需要同时评价任务结果和部署代价，指标应覆盖条件成功率、推理时延、回合级误报、报警提前量、人工接管和安全事件。

就本次检索的代表性样本而言，国外研究在动作分块模仿学习、移动操作全身遥操作、通用VLA模型与失败检测框架方面起步较早\cite{zhao2023act,fu2024mobilealoha,kim2024openvla,xu2025faildetect}；国内团队在双臂扩散基础模型、轻量与高效VLA结构和三维输入模仿学习方面形成了针对性较强的研究\cite{liu2025rdt,wen2025tinyvla,bu2025univla,liu2024robomamba,ze2024dp3,wang2024rise}，哈工大团队在精密装配示范学习方面也有持续积累\cite{zang2023promp}。上述比较限于所选代表性样本及相应的方法侧重。

将上述进展映射到本课题的目标场景，仍有三个问题有待验证。其一，在本文纳入比较的研究中，尚未发现把定位系统估计的工位停靠残差作为策略显式状态、同时在示范采集中系统覆盖随机停靠的方案，数据覆盖与显式条件的相对贡献也缺少对应消融。其二，按自由空间与精细阶段区别设置执行时域，并通过RTC处理流匹配动作块异步衔接，其成功率收益和计算代价需要在同一控制接口下评价。其三，RTC重叠区中新旧动作块一致性残差能否作为无需额外模型训练的失败信号，以及其回合级召回、误报和报警提前量，仍需在独立测试集上实证。

基于以上问题，本研究以已贯通并完成单回合功能验证的OpenPI \(\pi_{0.5}\)接口链路为基线，沿“工位停靠残差--精细阶段误差显现--错误动作持续”的传递链构建分层鲁棒增强框架，并在统一的残差注入协议、五子任务链和四台机器人条件下依次叠加验证。核心可检验问题分别落实为容忍边界\(R_{80}\)、精细子任务条件成功率与推理代价、检测召回率与回合误报率、报警提前量、恢复后整链成功率及安全事件数，所有改善结论均以预先规定的统计判据判断。

\section{主要内容及方案}

\subsection{研究主线与总体技术路线}

本课题的研究主线是工位停靠残差在\(\pi_{0.5}\)策略执行过程中的传递与显现。如图\ref{fig:error-chain-defense}所示，位置残差近似等幅改变目标相对位置，航向残差经目标距离转化为横向误差；当两者在进槽、落台等精细阶段超过任务允许量时，策略级开环执行又可能使错误动作持续至下一次重规划。针对三个环节，本课题构建分层鲁棒增强框架：停靠残差条件化模仿学习方法通过状态条件和示范分布提高策略对残差的适应能力；阶段自适应动作块控制方法通过调节执行时域并采用RTC连续衔接来缩短精细阶段的开环区间；预测一致性驱动的运行时容错方法通过监测RTC重叠区的跨块一致性残差触发模板化恢复。三项方法共享同一残差注入协议；压力测试档位用于采集成功和失败回合，检测器仅以其中独立的成功标定集确定阈值，失败回合只用于阈值冻结后的评价。

\begin{figure}[htbp]
  \centering
  \begin{tikzpicture}[
  font=\small,
  >={Stealth[length=2.2mm]},
  error/.style={draw=red!55!black,rounded corners=2pt,fill=red!4,
    minimum height=17mm,text width=34mm,align=center,inner sep=2.5pt},
  defense/.style={draw=blue!60!black,rounded corners=2pt,fill=blue!7,
    minimum height=25mm,text width=34mm,align=center,inner sep=2.5pt},
  validate/.style={draw=orange!75!black,rounded corners=2pt,fill=orange!10,
    minimum height=12mm,text width=118mm,align=center,inner sep=2.5pt},
  errarrow/.style={->,thick,red!65!black},
  block/.style={->,thick,blue!70!black},
  toval/.style={->,densely dashed,orange!80!black}
]
  \node[error] (e1) {\textbf{工位停靠残差}\\位置残差：厘米量级\\航向残差：度量级};
  \node[error,right=8mm of e1] (e2) {\textbf{精细阶段显现}\\位置残差近似等幅传递\\航向项随目标距离增加};
  \node[error,right=8mm of e2] (e3) {\textbf{错误动作持续}\\块内不接收新观测\\可能导致失败或损坏};
  \draw[errarrow] (e1) -- (e2);
  \draw[errarrow] (e2) -- (e3);

  \node[defense,below=10mm of e1] (d1) {\textbf{停靠残差条件化}\\\textbf{模仿学习方法}\\残差状态增广\\随机停靠示范};
  \node[defense,below=10mm of e2] (d2) {\textbf{阶段自适应}\\\textbf{动作块控制方法}\\执行时域调节\\RTC连续衔接};
  \node[defense,below=10mm of e3] (d3) {\textbf{预测一致性驱动的}\\\textbf{运行时容错方法}\\跨块残差检测\\模板恢复或停机};
  \draw[block] (d1) -- node[right,font=\footnotesize]{条件化适应} (e1);
  \draw[block] (d2) -- node[right,font=\footnotesize]{缩短开环} (e2);
  \draw[block] (d3) -- node[right,font=\footnotesize]{截断} (e3);

  \node[validate,below=8mm of d2] (val) {\textbf{整体验证}：C0基线 / C1条件化学习 / C2自适应控制 / C3检测与恢复，相同残差档位与五子任务链，四台机器人重复};
  \draw[toval] (d1.south) -- (d1.south |- val.north);
  \draw[toval] (d2) -- (val);
  \draw[toval] (d3.south) -- (d3.south |- val.north);
  \node[draw=red!50!black,densely dashed,rounded corners=3pt,fit=(e1)(e3),inner sep=2.5mm,
    label={[font=\footnotesize,text=red!60!black]above left:误差传递链}] {};
\end{tikzpicture}

  \caption{工位停靠残差传递链与分层鲁棒增强方法的对应关系}
  \label{fig:error-chain-defense}
\end{figure}

策略模型统一采用基于OpenPI微调的\(\pi_{0.5}\)，三项方法围绕其数据接口与流匹配动作块推理接口实现。所有方法均不改变策略网络主体结构：条件化学习方法扩展状态维度与归一化统计量，自适应控制方法和运行时容错方法作用于推理侧。主结果按“C0固定执行时域基线--C1停靠残差条件化学习--C2阶段自适应执行时域与RTC--C3检测与模板恢复”依次叠加；为分离组合内部因素，另设置短、中、长固定执行时域加RTC的C2-RS/RM/RL、阶段自适应加RTC的C2-A、离线只记录报警的C3-D和整链最多启用一次恢复的C3-R辅助配置。所有配置在相同残差档位、五子任务链和四台机器人上比较，结论分别限定为条件化学习、C2组合控制、检测器和恢复机制的对应效应。

\subsection{任务定义与工位停靠残差传递问题}
\label{subsec:task-definition}

PCBA上下料任务由四个物理工位上的五个操作子任务顺序组成，记子任务集合为\(\mathcal L=\{\text{上料车取料},\allowbreak\text{扫码},\allowbreak\text{检测台下料},\allowbreak\text{检测台上料},\allowbreak\text{下料车放料}\}\)，分别对应从上料百叶车夹取PCBA、扫码登记、从检测台取出上一块PCBA、将当前PCBA放入检测台以及将上一块PCBA放入下料百叶车，如图\ref{fig:pcba-station-flow}所示。外部任务状态机顺序给出语言条件\(\ell_t\in\mathcal L\)，\(\pi_{0.5}\)只生成低层连续动作，高层任务分解不纳入本课题评价。机器人导航至工位并锁定底盘后执行操作；检测台下料和上料在同一次停靠中连续完成，因而一条完整任务链包含四次停靠事件和五次子任务判定。

正式评价以一条完整的五子任务链为一个测试回合\(e\)。在线状态机和离线评价采用相互独立且在测试前冻结的信号。在线执行时，上料车取料与检测台下料在“夹爪闭合反馈进入冻结区间、腕部视觉中的PCBA完全离开原承载区域并持续3秒”后转移；扫码在扫码器返回与当前PCBA匹配的有效编码后转移；检测台上料以PCBA到达检测台机械限位并在释放后保持3秒为转移条件，下料车放料则以PCBA进入预先标定的目标槽位区域、达到冻结插入深度并在释放后保持3秒为转移条件。正式结果由不知道配置和残差档位的评价者依据同步视频、夹爪状态、扫码日志及工位传感记录复核同一组判据，不由在线状态机的转移结果直接充当成功标签。

每台机器人先用同一冻结C0检查点完成10个不计入正式统计的名义停靠回合，以各子任务时长中位数\(\widetilde T_m^{\mathrm{pilot}}\)的1.5倍作为超时上限\(T_m^{\max}\)，随后对所有配置保持不变。掉落、触发控制器接触保护或急停、人工介入以及超时均判为失败。每次残差注入最多允许三次导航尝试；三次后仍未进入预定档位则记为注入失败，不纳入策略成功率并按原随机化序列在该区组末尾补做。所有尝试及其参考残差均保留，用机器人、工位、档位和四停靠方向序列编号分层报告注入到位率，避免只保留易命中的停靠结果。

令\(Y_{e,m}\in\{0,1\}\)表示回合\(e\)的第\(m\)个子任务是否满足上述条件，则整链成功率和第\(m\)个子任务的条件成功率分别定义为
\begin{equation}
  \widehat S_{\mathrm{chain}}
  =
  \frac{1}{N}\sum_{e=1}^{N}\prod_{m=1}^{5}Y_{e,m},
  \qquad
  \widehat S_m^{\mathrm{cond}}
  =
  \frac{
    \sum_{e=1}^{N}Y_{e,m}\prod_{q=1}^{m-1}Y_{e,q}
  }{
    \sum_{e=1}^{N}\prod_{q=1}^{m-1}Y_{e,q}
  },
  \label{eq:task-success-metrics}
\end{equation}
其中空乘积取1；若某后序子任务的条件成功率分母为0，则报告为“不可估计”，不记为0。整链成功同时要求五个子任务均成功且全程无人工介入；导航到位率、注入失败率、人工接管率、急停次数、PCBA或夹具损伤数另行报告，使操作策略、实验施扰和系统安全三类结果不相互混淆。每个后序精细子任务在每个主比较档位至少保留30次有效到达；整链试验不足时，从同一冻结入口状态开展预先排定的子任务级补充试验，其结果只用于\(\widehat S_m^{\mathrm{cond}}\)，不并入整链成功率。

\begin{figure}[htbp]
  \centering
  \begin{tikzpicture}[
  font=\small,
  >={Stealth[length=2.2mm]},
  task/.style={draw=blue!60!black,rounded corners=2pt,fill=blue!7,
    minimum width=22mm,minimum height=13mm,text width=20mm,
    align=center,inner sep=2.5pt,outer sep=0pt},
  info/.style={rounded corners=2pt,minimum width=36mm,minimum height=13mm,
    text width=34mm,align=center,inner sep=2.5pt,outer sep=0pt},
  nominalbox/.style={info,draw=blue!60!black,fill=blue!7},
  actualbox/.style={info,draw=orange!75!black,fill=orange!10},
  residualbox/.style={info,draw=red!55!black,fill=red!4},
  linklabel/.style={font=\footnotesize,fill=white,inner sep=1pt,
    text=blue!65!black},
  mainarrow/.style={->,thick,blue!70!black},
  nominal/.style={->,thick,black},
  actual/.style={->,thick,densely dashed,orange!80!black},
  err/.style={->,thick,red!65!black}
]
  \node[task] (n1) {\textbf{上料百叶车}\\夹取PCBA};
  \node[task,right=8mm of n1] (n2) {\textbf{扫码工位}\\扫码登记};
  \node[task,right=8mm of n2] (n3) {\textbf{检测台}\\取出上一块};
  \node[task,right=8mm of n3] (n4) {\textbf{检测台}\\放入当前块};
  \node[task,right=8mm of n4] (n5) {\textbf{下料百叶车}\\放入上一块};
  \draw[mainarrow] (n1) -- node[linklabel,above]{导航} (n2);
  \draw[mainarrow] (n2) -- node[linklabel,above]{导航} (n3);
  \draw[mainarrow] (n3) -- node[linklabel,above]{原地} (n4);
  \draw[mainarrow] (n4) -- node[linklabel,above]{导航} (n5);

  % 工位停靠误差示意：以任一工位 k 为例
  \coordinate (o) at ($(n3.south)+(0,-30mm)$);
  \coordinate (b) at ($(o)+(12mm,-7mm)$);
  \draw[fill=gray!15,draw=gray!60]
    ($(o)+(-29mm,12mm)$) rectangle ($(o)+(29mm,23mm)$);
  \node[font=\footnotesize,align=center,text width=52mm]
    at ($(o)+(0,17.5mm)$)
    {第\(k\)个工位设备\\（百叶车 / 扫码枪 / 检测台）};
  \draw[nominal] (o) -- ++(0,10mm) node[left,font=\footnotesize]{\(x_{S_k}\)};
  \draw[nominal] (o) -- ++(-10mm,0) node[below,font=\footnotesize]{\(y_{S_k}\)};
  \fill (o) circle (0.5mm);
  \node[below left=0.5mm and 0mm of o,font=\footnotesize] {\(\{S_k\}\)};
  \begin{scope}[shift={(b)},rotate=15]
    \draw[actual] (0,0) -- (0,10mm) node[right,font=\footnotesize]{\(x_B\)};
    \draw[actual] (0,0) -- (-10mm,0) node[below,font=\footnotesize]{\(y_B\)};
    \draw[densely dashed,orange!80!black] (-7mm,-9mm) rectangle (7mm,2mm);
    \node[below,font=\footnotesize,text=orange!80!black,fill=white,
      inner sep=1pt] at (0,-11mm) {\(\{B\}\)实际底盘};
  \end{scope}
  \fill[orange!80!black] (b) circle (0.5mm);
  \coordinate (pturn) at (b |- o);
  \draw[err] (o) -| (b);
  \node[above=1mm,font=\scriptsize,text=red!65!black,fill=white,
    inner sep=1pt] at ($(o)!0.38!(pturn)$)
    {\(\delta\mathbf p_k^{\mathrm{ref}}\)};
  \draw[err] ($(b)+(90:7mm)$) arc[start angle=90,end angle=105,radius=7mm]
    node[pos=0.5,above,fill=white,inner sep=1pt,
      font=\scriptsize,text=red!65!black]{\(\delta\theta_k^{\mathrm{ref}}\)};

  \node[nominalbox,left=18mm of o,yshift=1mm] (nom)
    {\textbf{名义停靠位姿}\({}^{O}\mathbf T_{S_k}\)\\示范采集参考位姿};
  \node[actualbox,right=21mm of o,yshift=1mm] (act)
    {\textbf{实际参考位姿}\({}^{O}\mathbf T_B^{\mathrm{ref}}\)\\独立工位测量给出};
  \node[residualbox,below=4mm of nom] (res)
    {\textbf{三类停靠记录}\\导航指令\(\boldsymbol\delta_k^{\mathrm{cmd}}\)\\策略输入\(\widehat{\boldsymbol\delta}_k^{\mathrm{loc}}\)\\评价分档\(\boldsymbol\delta_k^{\mathrm{ref}}\)};
\end{tikzpicture}

  \caption{PCBA上下料五子任务链与工位停靠残差的定义}
  \label{fig:pcba-station-flow}
\end{figure}

为描述工位停靠残差，对第\(k\)个工位定义固连于环境的名义停靠坐标系\(\{S_k\}\)，该坐标系在示范采集时与机器人基座系重合，并不表示设备几何原点。记\({}^{A}\mathbf T_C\)为把\(\{C\}\)中坐标变换到\(\{A\}\)中的齐次变换；\({}^{O}\mathbf T_{S_k}\)给出名义停靠位姿，\({}^{O}\mathbf T_B^{\mathrm{ref}}\)由独立工位测量给出实际基座参考位姿，\({}^{O}\widehat{\mathbf T}_B^{\mathrm{loc}}\)由激光雷达定位系统给出其估计。由齐次变换复合关系，评价参考残差为
\begin{equation}
  \Delta\mathbf T_k^{\mathrm{ref}}
  =
  \left({}^{O}\mathbf T_{S_k}\right)^{-1}{}^{O}\mathbf T_B^{\mathrm{ref}}
  =
  \begin{bmatrix}
    \mathbf R(\delta\theta_k^{\mathrm{ref}}) & \delta\mathbf p_k^{\mathrm{ref}} \\
    \mathbf 0^{\mathrm T} & 1
  \end{bmatrix},
  \qquad
  \boldsymbol\delta_k^{\mathrm{ref}}
  =
    \left(\delta x_k^{\mathrm{ref}},\delta y_k^{\mathrm{ref}},\delta\theta_k^{\mathrm{ref}}\right)^{\mathrm T},
  \label{eq:station-frame-residual}
\end{equation}
式中，底盘停靠于水平地面，残差退化为平面位姿；\(\delta\mathbf p_k^{\mathrm{ref}}=(\delta x_k^{\mathrm{ref}},\delta y_k^{\mathrm{ref}})^{\mathrm T}\)为名义停靠系中的位置残差，\(\delta\theta_k^{\mathrm{ref}}\)为航向残差，\(\mathbf R(\cdot)\)为平面旋转矩阵。以参考测量的平面位姿\((x_B^{\mathrm{ref}},y_B^{\mathrm{ref}},\theta_B^{\mathrm{ref}})\)和名义位姿\((x_{S_k},y_{S_k},\theta_{S_k})\)表示，并用\(\operatorname{wrap}_{[-\pi,\pi)}\)取航向差主值时，
\begin{equation}
  \delta\theta_k^{\mathrm{ref}}=\operatorname{wrap}_{[-\pi,\pi)}(\theta_B^{\mathrm{ref}}-\theta_{S_k}),
  \qquad
  \delta\mathbf p_k^{\mathrm{ref}}
  =
  \mathbf R(\theta_{S_k})^{\mathrm T}
  \left(
    \begin{bmatrix}x_B^{\mathrm{ref}}\\y_B^{\mathrm{ref}}\end{bmatrix}
    -
    \begin{bmatrix}x_{S_k}\\y_{S_k}\end{bmatrix}
  \right).
  \label{eq:residual-planar}
\end{equation}

为区分实验干预、策略可见量和评价参考量，每次停靠同时记录导航目标注入量\(\boldsymbol\delta_k^{\mathrm{cmd}}\)、由\({}^{O}\widehat{\mathbf T}_B^{\mathrm{loc}}\)按同一变换计算的停靠残差定位估计\(\widehat{\boldsymbol\delta}_k^{\mathrm{loc}}\)，以及上述参考残差\(\boldsymbol\delta_k^{\mathrm{ref}}\)。正式实验前使用经校准量具或等效外部测量装置校核独立测量通道，要求位置与航向的扩展不确定度分别不大于2毫米和\(0.2^\circ\)（覆盖因子\(k=2\)），10次重复测量的标准差分别不大于1毫米和\(0.1^\circ\)。正式测试按\(\boldsymbol\delta_k^{\mathrm{ref}}\)而非导航命令归档：目标平移档位的轴向幅值接纳半宽为3毫米、正交残差不超过3毫米且航向残差不超过\(0.3^\circ\)；目标航向档位的接纳半宽为\(0.3^\circ\)，平移残差不超过3毫米。定位估计相对于参考残差的绝对误差中位数和95分位数随主结果报告，从而不把导航命令或定位估计直接视为物理真值。

工位停靠残差对操作精度的影响可由目标点在基座系中的位置变化刻画。以下为简化符号，令\(\delta\mathbf p_k=\delta\mathbf p_k^{\mathrm{ref}}\)、\(\delta\theta_k=\delta\theta_k^{\mathrm{ref}}\)。设槽位、检测台卡口等操作目标固连于工位，在名义停靠系中的平面位置为\(\mathbf p^{S_k}_G\)。示范策略隐含地假定基座与名义停靠系重合，其动作以\(\mathbf p^{S_k}_G\)为目标；实际停靠后，同一目标在基座系中的位置为\(\mathbf R(\delta\theta_k)^{\mathrm T}(\mathbf p^{S_k}_G-\delta\mathbf p_k)\)，与策略隐含目标之差为
\begin{equation}
  \mathbf e_p
  =
  \mathbf R(\delta\theta_k)^{\mathrm T}
  \left(\mathbf p^{S_k}_G-\delta\mathbf p_k\right)
  -\mathbf p^{S_k}_G
  =
  \left(\mathbf R(\delta\theta_k)^{\mathrm T}-\mathbf I\right)
  \mathbf p^{S_k}_G
  -
  \mathbf R(\delta\theta_k)^{\mathrm T}\delta\mathbf p_k,
  \label{eq:residual-ee-error}
\end{equation}
由\(\|\mathbf R(\delta\theta)^{\mathrm T}-\mathbf I\|_2=2|\sin(\delta\theta/2)|\)可得严格上界
\begin{equation}
  \left\|\mathbf e_p\right\|
  \le
  \left\|\delta\mathbf p_k\right\|
  +
  2\left|\sin\frac{\delta\theta_k}{2}\right|
  \left\|\mathbf p^{S_k}_G\right\|
  \le
  \left\|\delta\mathbf p_k\right\|
  +
  \left|\delta\theta_k\right|\,\left\|\mathbf p^{S_k}_G\right\|,
  \label{eq:residual-error-bound}
\end{equation}
式\eqref{eq:residual-error-bound}表明，位置残差近似等幅传递，航向残差项随目标距离增加。例如，当目标距离为0.5米时，\(0.5^\circ\)航向残差对应的该项上界约为4.4毫米。各工位的实际目标距离、PCBA与槽位有效间隙在正式测试前测量并写入冻结的评价协议，实验据此判断几何误差是否进入精细操作敏感区。上述关系只描述水平面停靠残差的几何效应；测试中高度、横滚、俯仰及工位标定状态保持不变，不将其解释为全部末端误差来源。

\subsection{停靠残差条件化模仿学习方法}
\label{subsec:improvement-one}

条件化学习方法从数据分布与状态条件两个方面处理残差。数据层面，随机停靠示范对每次停靠事件独立采样
\begin{equation}
  \delta_{x,k}^{\mathrm{cmd}},\delta_{y,k}^{\mathrm{cmd}}
  \sim\mathcal U(-3\,\mathrm{cm},3\,\mathrm{cm}),
  \qquad
  \delta_{\theta,k}^{\mathrm{cmd}}
  \sim\mathcal U(-3^\circ,3^\circ),
  \label{eq:training-residual-distribution}
\end{equation}
即\(\pm3\)厘米是两个位置分量各自的范围，平移幅值最大为\(3\sqrt2\)厘米。每个示范回合对应一个低层子任务并关联其所在停靠事件，实际参考残差、定位估计残差和导航指令量均随回合记录；操作者在遥操作中对残差的修正同步成为动作标签。每个子任务至少采集50个有效示范回合，并在位置四个象限与航向正负方向间均衡；固定停靠组取相同回合数，以避免数据量差异与残差覆盖效应混杂。

状态层面，将残差拼接到已贯通OpenPI接口的14维本体状态，形成残差增广状态与观测
\begin{equation}
  \begin{aligned}
    \tilde{\mathbf x}_t
    &=
    \left[\mathbf x_t^{\mathrm T},\bar{\boldsymbol\delta}_k^{\mathrm T}\right]^{\mathrm T}
    \in\mathbb R^{17},
    \qquad
    \bar{\boldsymbol\delta}_k
    =
    \left(
      \frac{\widehat\delta^{\mathrm{loc}}_{x,k}-\mu_x}{\sigma_x^\star},
      \frac{\widehat\delta^{\mathrm{loc}}_{y,k}-\mu_y}{\sigma_y^\star},
      \frac{\widehat\delta^{\mathrm{loc}}_{\theta,k}-\mu_\theta}{\sigma_\theta^\star}
    \right)^{\mathrm T},\\
    \tilde{\mathbf o}_t
    &=
    \left\{
      I^{\mathrm{head}}_t,
      I^{\mathrm{wrist},L}_t,
      I^{\mathrm{wrist},R}_t,
      \tilde{\mathbf x}_t,
      \ell_t
    \right\},
  \end{aligned}
  \label{eq:residual-augmented-state}
\end{equation}
式中，\(\mathbf x_t\)为原14维双臂末端位姿与夹爪状态；\(\mu_x,\mu_y,\mu_\theta\)和各标准差只由训练划分估计，\(\sigma_x^\star=\max(\sigma_x,2\,\mathrm{mm})\)、\(\sigma_y^\star=\max(\sigma_y,2\,\mathrm{mm})\)、\(\sigma_\theta^\star=\max(\sigma_\theta,0.2^\circ)\)，以避免近零方差造成数值放大。相同统计量原样用于验证、标定与测试。定位系统估计的停靠残差在一次停靠期间保持不变，仅在机器人重新导航并停靠后更新；检测台连续执行的下料与上料共用同一估计。状态接口扩展为17维，策略网络主体和既有流匹配训练目标保持不变；视觉输入能否隐式补偿残差以及显式状态是否提供增量收益均由消融结果判断。

条件化学习方法的实验包括平移容忍曲线、航向容忍曲线和消融。平移曲线取参考残差幅值\(\rho=\|\delta\mathbf p_k^{\mathrm{ref}}\|\in\{0,1,2,3,4,5\}\)厘米，非零档位在名义停靠系的\(\pm x\)、\(\pm y\)四个方向均衡分配；航向曲线取\(|\delta\theta_k^{\mathrm{ref}}|\in\{0,1,2,3,4,5\}^\circ\)，正负方向均衡。两条曲线共用零残差档位；零档要求平移与航向均满足第\ref{subsec:task-definition}节的非施扰轴容差。一条正式测试链的四次停靠均采用同一幅值档位，各停靠方向按预先生成且编号的均衡序列分配，统计分层和匹配均使用该“四停靠方向序列编号”，不把一条链归入某个单一方向。

C0、C1与C2-A三种主曲线配置均覆盖三个配对训练随机种子；在每个“机器人\(\times\)配置\(\times\)档位”单元内，每个种子至少完成4条有效链，即三种子合计不少于12条、四台机器人合计不少于48条。每个机器人和种子另完成8条自然导航停靠链，即每个配置合计96条；自然链只作为耦合残差补充测试，不用于计算单轴容忍边界。配置、档位、种子与方向序列按机器人和日期分区组随机化，各种子在汇总统计中等权。消融(b)仅在共用零档、3厘米与\(3^\circ\)两个训练边界档位采用同样分配，以回答预先规定的数据覆盖问题。

对平移或航向等级统一记为\(r\)，相应整链成功率记为\(\widehat S(r)\)。主验收的\(\widehat S(r)\)按四台机器人和三个种子等权汇总，源机器人和三台未参与训练的目标机器人组另行报告。容忍边界与归一化梯形曲线面积定义为
\begin{equation}
  \begin{aligned}
    R_{80}
    &=
    \max\left\{
      r:
      \operatorname{LCB}^{\mathrm{boot}}_{0.95}\!\left[\widehat S(r')\right]\ge0.80,
      \ \forall r'\le r
    \right\},\\
    \operatorname{AUC}_{\mathrm{rob}}
    &=
    \frac{1}{5}
    \left[
      \frac{\widehat S(0)+\widehat S(5)}{2}
      +
      \sum_{r=1}^{4}\widehat S(r)
    \right].
  \end{aligned}
  \label{eq:robustness-metrics}
\end{equation}
其中\(R_{80}\)分别以厘米和度报告；若零档也不满足条件，则记为“未达到”，不把空集解释为0。\(\operatorname{LCB}^{\mathrm{boot}}_{0.95}\)与配置差区间均由10000次分层簇自举获得：先按机器人和训练种子抽取簇，再在日期与四停靠方向序列内抽取回合，以保持实验的聚类结构。消融按表\ref{tab:residual-ablation}设置三组配置：(a)即C0，(c)即C1，(b)为识别数据覆盖贡献的辅助配置；(a)与(b)比较数据覆盖，(b)与(c)比较显式条件。三组均从同一OpenPI预训练检查点独立微调三个随机种子，示范回合数、训练步数及其余超参数一致，且(b)与(c)使用完全相同的随机停靠回合。

\begin{longtable}{@{}p{1.4cm}p{3.8cm}p{2.4cm}p{5.8cm}@{}}
  \caption{停靠残差条件化模仿学习的消融配置}\label{tab:residual-ablation}\\
  \toprule
  配置 & 训练数据停车方式 & 状态含残差 & 考察问题 \\
  \midrule
  \endfirsthead
  \multicolumn{4}{c}{续表~\thetable}\\
  \toprule
  配置 & 训练数据停车方式 & 状态含残差 & 考察问题 \\
  \midrule
  \endhead
  (a) & 固定名义停靠 & 否 & 现有流水线基线，给出名义条件下的性能与残差注入下的退化速度 \\
  (b) & \(\delta x,\delta y\in[-3,3]\)厘米，\(\delta\theta\in[-3,3]^\circ\) & 否 & 仅靠数据覆盖，策略能否从图像中隐式推断并补偿残差 \\
  (c) & 与(b)相同的随机停靠回合 & 是 & 显式残差输入在数据覆盖之外能否提高\(R_{80}\)与归一化曲线面积 \\
  \bottomrule
\end{longtable}

\subsection{阶段自适应动作块控制方法}
\label{subsec:improvement-two}

自适应控制方法不改变训练，只在推理侧按阶段调整每次推理后的执行时域。将时刻\(t\)的执行阶段记为\(\phi_t\in\{\mathrm{free},\mathrm{fine}\}\)，其中精细阶段对应进槽、出槽、落台等对准操作，自由阶段对应工位内的接近、搬运与撤离。含精细区段的子任务集合固定为\(\mathcal L_{\mathrm{fine}}=\{\text{上料车取料},\text{检测台下料},\text{检测台上料},\text{下料车放料}\}\)。为避免距离阈值附近反复切换，阶段判定采用进入阈值\(d_{\mathrm{enter}}=10\)厘米和退出阈值\(d_{\mathrm{exit}}=12\)厘米的滞回规则：
\begin{equation}
  \phi_t
  =
  \begin{cases}
    \mathrm{fine}, &
    \ell_t\in\mathcal L_{\mathrm{fine}}
    \ \text{且}\
    \left\|{}^{B}\mathbf p_{E_\alpha,t}-{}^{B}\mathbf p_{G_{\ell_t,\alpha}}\right\|\le d_{\mathrm{enter}},\\
    \mathrm{free}, &
    \ell_t\notin\mathcal L_{\mathrm{fine}}
    \ \text{或}\
    \left\|{}^{B}\mathbf p_{E_\alpha,t}-{}^{B}\mathbf p_{G_{\ell_t,\alpha}}\right\|\ge d_{\mathrm{exit}},\\
    \phi_{t-1}, & \text{其他},
  \end{cases}
  \qquad
  s_t=s_{\phi_t},
  \quad
  s_{\mathrm{fine}}<s_{\mathrm{free}}\le H,
  \label{eq:phase-rule}
\end{equation}
式中，\({}^{B}\mathbf p_{E_\alpha,t}\in\mathbb R^3\)为当前操作臂\(\alpha\in\{L,R\}\)在基座系中的末端位置，\({}^{B}\mathbf p_{G_{\ell_t,\alpha}}\in\mathbb R^3\)为由工位标定和定位系统估计的停靠残差换算得到、随子任务与操作臂变化的三维目标位置；第\ref{subsec:task-definition}节的\(\mathbf p_G^{S_k}\)仅是该目标在水平面的投影，用于二维误差传播分析。\(H\)为模型每次生成的预测时域，\(s_{\mathrm{free}}\)与\(s_{\mathrm{fine}}\)分别为两类阶段的执行时域。每个子任务初始阶段置为\(\mathrm{free}\)，阶段信号由外部任务状态机和距离规则共同确定，不依赖原始\(\pi_{0.5}\)的高层子任务预测。10厘米与12厘米阈值在正式测试前通过独立预试验检查工作空间和误触发情况，此后不随测试结果调整。

设第\(i\)次推理在控制步索引\(n_i\)启动，对应墙钟时刻为\(t_i=n_i\Delta t\)，得到长度为\(H\)的动作块\(\widehat{\mathbf A}_{n_i}\)；下一次推理索引为\(n_{i+1}=n_i+s_{\phi_{n_i}}\)。将该次推理耗时折算为控制步数
\begin{equation}
  d_i
  =
  \left\lceil\frac{\tau_{\mathrm{inf},i}}{\Delta t}\right\rceil,
  \qquad
  T_\phi=s_\phi\Delta t,
  \qquad
  d_i\le s_\phi,
  \quad
  s_\phi+d_i\le H,
  \label{eq:replan-period}
\end{equation}
式中，\(\Delta t\)为控制周期，\(\tau_{\mathrm{inf},i}\)从推理请求发出计至动作块进入执行队列，包含图像预处理、通信、模型推理、RTC引导和排队，\(d_i\)为相应控制步延迟，\(T_\phi\)为重规划周期。正式实验前在目标硬件上采集不少于1000次稳态推理，令
\begin{equation}
  d_\star
  =
  \left\lceil\frac{Q_{0.99}(\tau_{\mathrm{inf}})}{\Delta t}\right\rceil,
  \quad
  s_{\mathrm{short}}=d_\star,
  \quad
  s_{\mathrm{long}}=H-d_\star,
  \quad
  s_{\mathrm{middle}}
  =
  \left\lfloor\frac{s_{\mathrm{short}}+s_{\mathrm{long}}}{2}\right\rfloor,
  \label{eq:execution-horizon-selection}
\end{equation}
并在冻结参数前显式检查\(s_{\mathrm{short}}<s_{\mathrm{middle}}<s_{\mathrm{long}}\)，即\(d_\star<\lfloor H/2\rfloor<H-d_\star\)；若不满足，则先优化推理链路并重新测量，不进入正式测试。分阶段配置固定取\(s_{\mathrm{fine}}=s_{\mathrm{short}}\)、\(s_{\mathrm{free}}=s_{\mathrm{long}}\)。所有时域在最终测试前冻结；运行中若单次推理越过式\eqref{eq:replan-period}的期限，则先执行旧块尚有效的尾部并记录一次截止期违约。若新块返回前旧块尾部耗尽，控制器立即按冻结减速参数进入位姿保持，当前子任务记为失败，并在未决请求返回或被取消前抑制新的重叠请求，不把未定义动作发送给机器人。

执行采用适用于流匹配VLA的RTC机制\cite{black2025rtc}：新一轮推理进行时，机器人继续执行前一动作块；在生成新块时，将延迟\(d_i\)内已经确定的动作冻结为前缀，对其后的重叠区施加逐步减弱的引导，再由\(\pi_{0.5}\)流匹配动作头补全后续动作。该机制不对两个独立流匹配样本作直接平均。图\ref{fig:phase-chunk-timing}给出了两类阶段的执行时序：自由阶段采用较长执行时域以减少推理次数，精细阶段采用较短执行时域，使新观测更快进入控制回路，并由RTC连续衔接重叠动作块。

\begin{figure}[htbp]
  \centering
  \begin{tikzpicture}[
  font=\small,
  >={Stealth[length=2.2mm]},
  exec/.style={fill=blue!55!black,draw=blue!60!black},
  pend/.style={fill=blue!12,draw=blue!60!black},
  infer/.style={->,thick,orange!80!black},
  note/.style={font=\footnotesize,align=center}
]
  % 阶段区域
  \fill[gray!10] (0.2,-0.5) rectangle (5.3,3.55);
  \fill[red!5] (5.3,-0.5) rectangle (12.4,3.55);
  \draw[densely dashed,gray!60!black] (5.3,-0.5) -- (5.3,3.55);
  \node[note,anchor=south west,text=gray!60!black] at (0.3,3.65)
    {自由阶段\(\phi_t=\mathrm{free}\)：块长\(H\)，\\执行\(h_{\mathrm{free}}\approx H\)步（深色）后再推理};
  \node[note,anchor=south west,text=red!60!black] at (5.45,3.65)
    {精细阶段\(\phi_t=\mathrm{fine}\)：每执行\(h_{\mathrm{fine}}\)步即重推理，\\未执行部分（浅色）与后续块重叠};
  % 时间轴
  \draw[->,thick] (0,0) -- (12.8,0) node[below left,font=\footnotesize]{\(t\)};
  \foreach \x in {0.4,1.1,...,12.3} \draw (\x,0.06) -- (\x,-0.06);
  % 自由阶段：一条长块
  \draw[pend] (0.4,0.5) rectangle (5.6,0.85);
  \draw[exec] (0.4,0.5) rectangle (5.3,0.85);
  \draw[infer] (0.4,-0.45) -- (0.4,-0.05);
  % 精细阶段：错位重叠的短块，H=2.8，h_fine=0.7
  \foreach \s/\y in {5.3/0.5,6.0/1.0,6.7/1.5,7.4/2.0,8.1/2.5,8.8/3.0} {
    \draw[pend] (\s,\y) rectangle (\s+2.8,\y+0.35);
    \draw[exec] (\s,\y) rectangle (\s+0.7,\y+0.35);
    \draw[infer] (\s,-0.45) -- (\s,-0.05);
  }
  % 当前时刻 t 处的重叠块集合
  \draw[densely dashed,gray!60!black] (8.5,-0.5) -- (8.5,3.45);
  \draw[draw=gray!60!black,densely dotted,rounded corners=2pt] (8.3,0.95) rectangle (8.7,2.9);
  \node[note,anchor=south east] at (8.25,2.95) {\(\mathcal C_t\)};
  \node[note,anchor=north] at (8.5,-0.55) {当前时刻\(t\)};
  \node[note,anchor=north] at (2.8,-0.55) {橙色箭头：推理时刻\(t_i\)，\(t_{i+1}=t_i+h_{\phi_{t_i}}\)};
  \node[note,anchor=north,text width=34mm] at (11.0,-0.55) {\(\mathcal C_t\)：覆盖\(t\)的重叠块集合，\(|\mathcal C_t|=\lceil H/h_{\mathrm{fine}}\rceil\)，供改进点三计算不一致度};
\end{tikzpicture}

  \caption{分阶段执行时域与RTC跨动作块连续衔接时序}
  \label{fig:phase-chunk-timing}
\end{figure}

自适应控制方法在每个C1训练种子对应的同一检查点上比较四种RTC执行配置：C2-RS、C2-RM和C2-RL分别在全任务采用式\eqref{eq:execution-horizon-selection}的短、中、长固定时域，三者统称C2-R；C2-A按式\eqref{eq:phase-rule}在短、长时域间切换。四组均启用相同RTC，因而C2-A与固定组的差异只归因于阶段切换；C1与C2的主配置差异解释为“RTC与阶段自适应组成的控制模块”总体效应，不单独归因于RTC。另以同一批已记录观测在\(s=s_{\mathrm{middle}}\)条件下进行RTC开关配对回放，测量RTC本身的服务开销。

C2-A按前述完整曲线样本量测试；三种固定时域对照只在共用零档、3厘米、5厘米、\(3^\circ\)和\(5^\circ\)五个代表档位测试。每个“机器人\(\times\)固定配置\(\times\)档位\(\times\)训练种子”单元至少4条有效链，即每个固定配置至少240条链。C2-A与C2-RS、C2-RL在五档等权汇总的精细子任务条件成功率上构成两项主要比较，C2-RM只用于描述时域--性能趋势，不进入主要验收。

稳态推理频率、首动作延迟、单次推理端到端服务时延、每链累计服务时间与GPU执行时间定义为
\begin{equation}
  \begin{aligned}
    f_{\mathrm{inf}}
    &=
    \frac{N_{\mathrm{inf}}-1}{t_{\mathrm{last}}-t_{\mathrm{first}}},
    &
    \tau_{\mathrm{first}}
    &=
    t_{\mathrm{first\ command}}-t_{\mathrm{subtask\ request}},\\
    C_{\mathrm{srv}}^e
    &=
    \sum_{i\in e}\tau_{\mathrm{inf},i},
    &
    C_{\mathrm{GPU}}^e
    &=
    \sum_{i\in e}\tau_{\mathrm{GPU},i}.
  \end{aligned}
  \label{eq:runtime-metrics}
\end{equation}
式中全部时间量以秒计算，表格展示时可换算为毫秒；频率以赫兹报告，时延、累计服务时间和GPU执行时间报告中位数与95分位数，同时报告推理调用数及截止期违约率。任务结果使用式\eqref{eq:task-success-metrics}，并在C2-A上重做残差容忍曲线。两项主要成功率比较给出95\%置信区间并采用Holm校正；每台机器人、每个配置还记录不少于1000次稳态推理，RTC开关基准按同一请求配对，避免把同一回合内的推理时刻误作成功率的独立样本。

\subsection{预测一致性驱动的运行时容错方法}
\label{subsec:improvement-three}

运行时容错方法利用RTC连续执行产生的重叠动作块构造失败信号，而不直接平均不同的流匹配样本。设第\(i-1\)次RTC输出动作块在执行\(s_{i-1}=s_{\phi_{n_{i-1}}}\)步后留下尾部\(\widehat{\mathbf A}^{\mathrm{prev}}_i\)，第\(i\)次推理在固定参考噪声\(\boldsymbol\epsilon_{\mathrm{ref}}\)下、RTC引导前得到原始动作提议\(\widehat{\mathbf A}^{\mathrm{raw}}_i\)。两者均按绝对控制步\(n_i+j\)对齐，重叠索引集合及逐步残差定义为
\begin{equation}
  \begin{aligned}
    \mathcal J_i
    &=
    \left\{j:0\le j<H-s_{i-1}\right\},\\
    \mathbf e_{i,j}
    &=
    \operatorname{diff}\!\left(
      \widehat{\mathbf a}^{\mathrm{raw}}_{i,j},
      \widehat{\mathbf a}^{\mathrm{prev}}_{i,j}
    \right),
    \qquad j\in\mathcal J_i.
  \end{aligned}
  \label{eq:cross-chunk-residual}
\end{equation}
算子\(\operatorname{diff}\)对位置分量作普通差，对滚转、俯仰和航向分量作\([-\pi,\pi)\)内的包角差，对夹爪分量按开合状态不一致记0或1，从而避免欧拉角跨界和离散夹爪状态造成伪残差。固定\(\boldsymbol\epsilon_{\mathrm{ref}}\)用于抑制流匹配采样随机性对检测分数的影响。C3-D使用冻结C2-A轨迹中不可更改的图像、状态、前一RTC块、检查点和时间戳，在目标硬件上顺序回放原始分支并离线计算报警，因而检测评价不改变被评价动作；C3-R在线运行时，原始分支与RTC执行分支放入同一批请求，实际执行只采用RTC引导后的动作块，批处理开销进入C3-R的实测时延。RTC掩码权重\(\omega_{i,j}\in[0,1]\)在延迟内的冻结前缀取1，在可编辑重叠区逐步衰减，并在重叠区外取0\cite{black2025rtc}。因此，该权重既规定新旧动作的连续约束，也确定一致性检测重点关注的近期动作。

跨动作块不一致度定义为重叠区归一化残差的加权均值，并固定在连续\(L=3\)次推理上平滑：
\begin{equation}
  r_i
  =
  \frac{
    \sum_{j\in\mathcal J_i}\omega_{i,j}
    \left\|\boldsymbol\Lambda^{-1/2}\mathbf e_{i,j}\right\|_2^2
  }{
    \sum_{j\in\mathcal J_i}\omega_{i,j}
  },
  \qquad
  \widetilde r_i
  =
  \frac{1}{L}
  \sum_{q=i-L+1}^{i}r_q.
  \label{eq:cross-chunk-inconsistency}
\end{equation}
式中，连续动作维度取\(\Lambda_{qq}=\max(\widehat\sigma_q^2,\varepsilon_q^2)\)，\(\widehat\sigma_q^2\)只由训练集估计，\(\varepsilon_q\)取对应控制接口的最小分辨率；夹爪维度采用固定单位权重。由于\(L=3\)且报警需要两个连续平滑分数，首次可判定报警从第4次推理开始。相邻观测条件下对同一未来时段的预测可能在正常状态下较接近、在分布偏移或卡滞前增大；该分数只是待实验验证的失效先兆，不预设其对所有失败类型均成立。

阈值只由成功标定回合确定，不使用失败标签。对机器人\(b\)单独划出与最终测试日期不重叠的成功标定集\(\mathcal E_{\mathrm{cal,succ}}^b\)，每台机器人不少于40条任务链。固定连续超阈次数\(n_{\mathrm{hold}}=2\)，并定义两次连续高分与全回合最大统计量
\begin{equation}
  \begin{aligned}
    g_i^e
    &=
    \min\!\left(\widetilde r_{i-1}^e,\widetilde r_i^e\right),\\
    M_b^e
    &=
    \max_i g_i^e,
    \qquad e\in\mathcal E_{\mathrm{cal,succ}}^b,\\
    \eta_b
    &=
    Q_{0.95}\!\left(
      \left\{M_b^e:e\in\mathcal E_{\mathrm{cal,succ}}^b\right\}
    \right),\\
    \text{报警}
    &\iff
    g_i>\eta_b.
  \end{aligned}
  \label{eq:threshold-p95}
\end{equation}
式\eqref{eq:threshold-p95}直接对所有阶段的全回合最大值取阈值，避免多个阶段阈值的报警并集突破回合预算；\(Q_{0.95}\)固定采用最近秩次序统计量，因而成功标定集内严格超阈回合的比例不超过5\%。\(L\)、\(n_{\mathrm{hold}}\)和所有阈值在查看失败测试结果前冻结。按机器人标定属于部署校准而非策略微调，主结果同时报告每台机器人40条成功链的校准成本；另以四台机器人等权合并的统一阈值作敏感性分析，不据测试结果在两种阈值间择优。

检测质量在C3-D离线影子模式下评价：冻结检测器顺序回放独立C2-A轨迹，只产生带时间戳的报警而不改变动作，因而原始失败时刻仍可观察。对第\(i\)个回放请求，报警可用时刻固定为原C2-A请求时刻加该请求在目标硬件上实测的批处理服务时延\(\tau_i^{\mathrm{batch}}\)，而不是采用离线作业的排队墙钟；由此在不改变原轨迹的条件下计入检测计算延迟。阈值标定集与检测测试集按日期、回合和四停靠方向序列完全隔离；独立测试集汇总不少于120个成功回合和80个残差诱发的精细操作失败回合，每台机器人至少包含30个成功回合和20个失败回合。导航、通信、扫码器或硬件故障不并入主检测集合，而作为其他故障单列。压力测试档位在试验前确定并记录全部回合，不按观察到的成败事后筛选。

每条链在首次报警后不重新布防，C3-D只记录一次报警事件，C3-R则由该事件进入一次恢复状态机。设\(t_f^e\)为失败起始时刻，\(t_a^e\)为该回合第一次报警结果实际可供控制器响应的时刻，无报警时取\(+\infty\)。两名不知道报警分数和配置的标注者依据同步视频与控制日志独立确定\(t_f^e\)：PCBA开始挤压或掉落、末端跟踪误差连续3个控制周期超过冻结的保护阈值、或物理状态已无法在当前子任务内恢复时，取最早发生者；分歧超过一个控制周期时由第三人裁定。停止试验覆盖最大代表载荷、最高测试速度与四类精细阶段姿态，每种组合不少于10次、合计不少于40次；将报警可用到末端速度降至限值5\%以内的最大实测时间再加一个控制周期定义为保守停止界\(T_{\mathrm{stop,safe}}\)。预警窗口固定为\(W_{\mathrm{pre}}=2\)秒，及时检出要求
\begin{equation}
  t_f^e-W_{\mathrm{pre}}
  \le
  t_a^e
  \le
  t_f^e-T_{\mathrm{stop,safe}}.
  \label{eq:timely-alarm-window}
\end{equation}
监测有效时段从首次具备报警判定条件起，到成功链完成或失败链的\(t_f^e\)止。失败回合中早于预警窗口或晚于安全停止界的首次报警、以及成功回合中的首次报警均定义为非窗口报警；失败发生后的信号不计入报警数或监测时长。由于每链最多记录一次报警，非窗口报警不会被连续超阈重复累计。回合级召回率、成功回合误报率、单位监测小时非窗口报警次数和逐真阳性提前量定义为
\begin{equation}
  \begin{aligned}
    \mathrm{Rec}_{\mathrm{ep}}
    &=
    \frac{\left|\{e\in\mathcal E_{\mathrm{test,fail}}:\text{满足式\eqref{eq:timely-alarm-window}}\}\right|}
         {\left|\mathcal E_{\mathrm{test,fail}}\right|},\\
    \mathrm{FPR}_{\mathrm{ep}}
    &=
    \frac{\left|\{e\in\mathcal E_{\mathrm{test,succ}}:\text{至少报警一次}\}\right|}
         {\left|\mathcal E_{\mathrm{test,succ}}\right|},\\
    \mathrm{FAR}_{\mathrm h}
    &=
    \frac{N_{\mathrm{nonwindow}}}{T_{\mathrm{monitor}}}\quad(\mathrm h^{-1}),
    \qquad
    T_{\mathrm{lead}}^e=t_f^e-t_a^e\quad(\mathrm s).
  \end{aligned}
  \label{eq:detector-metrics}
\end{equation}
式中\(T_{\mathrm{monitor}}\)为上述所有测试回合有效监测时段之和并换算为小时。二项指标报告Wilson 95\%置信区间；\(T_{\mathrm{lead}}^e\)只在及时检出的失败回合上报告中位数、四分位距和5分位数。

检测参数冻结后，在独立的C3-R回合中启用固定模板恢复，并与按机器人、模型检查点及训练种子、工位、残差档位、四停靠方向序列和日期区组匹配的C2-A无恢复回合随机交错执行。报警后先停止当前动作块；仅当夹爪状态正常、回退路径经离线检查且未发生不可逆卡滞时，沿工位竖直方向以不超过50毫米/秒退至安全高度\(z_{\mathrm{safe}}\)，该高度取工位最高障碍物上方50毫米；随后回到冻结的子任务入口位姿，重新获取观测并重试当前子任务。每条五子任务链最多允许一次恢复尝试；恢复成功后继续后序子任务但不再允许第二次恢复，后续任何报警均立即停机并交由人工处理。若初次报警时安全前提不满足、恢复超过由最大代表载荷下10次演练确定的\(1.5\)倍中位时间上限，或重试再次报警，也立即停机并交由人工处理。

C3-R在共用零档、3厘米、5厘米、\(3^\circ\)和\(5^\circ\)五个代表档位完成不少于240对匹配链，即每个“机器人\(\times\)档位\(\times\)训练种子”单元不少于4对；其中实际触发且满足恢复安全前提的恢复尝试总数不少于60次、每台机器人不少于12次，否则条件恢复成功率只作描述而不进入验收。最终样本量在不查看正式结果的前提下，按双侧\(\alpha=0.05\)、功效0.80、主要效应与非劣界及预试验基线率计算；功效计算大于上述下限时采用较大值。

条件恢复成功率与恢复后的最终整链成功率定义为
\begin{equation}
  S_{\mathrm{recover}}
  =
  \frac{N_{\mathrm{successful\ retries}}}{N_{\mathrm{recovery\ attempts}}},
  \qquad
  S_{\mathrm{final}}
  =
  \frac{N_{\mathrm{chains\ completed\ with\ at\ most\ one\ retry}}}
       {N_{\mathrm{all\ C3\ test\ chains}}}.
  \label{eq:recovery-metrics}
\end{equation}
同时报告首次尝试成功率、人工接管率、恢复诱发碰撞或急停次数以及PCBA或夹具损伤数，防止额外重试掩盖原策略退化。对一条有效链，\(T_{\mathrm{chain}}\)从首次工位导航请求发出计至第五个子任务的3秒稳定判定结束，包含四次导航、扫码等待、报警停止和自动恢复，不包含实验复位或人工处理时间；未触发恢复与触发恢复的回合分别汇总。令\(j\)索引同一输入在目标硬件上的配对基准请求，\(p(e)\)为C3-R回合\(e\)的预先匹配C2-A回合，则RTC、监测和恢复开销定义为配对差的95分位数
\begin{equation}
  \begin{aligned}
    \Delta\tau_{\mathrm{RTC},95}
    &=Q_{0.95}\!\left(
      \left\{\tau_j^{\mathrm{RTC,on}}-\tau_j^{\mathrm{RTC,off}}\right\}_{j=1}^{J}
    \right),\\
    \Delta\tau_{\mathrm{mon},95}
    &=Q_{0.95}\!\left(
      \left\{\tau_j^{\mathrm{batch}}-\tau_j^{\mathrm{RTC}}\right\}_{j=1}^{J}
    \right),\\
    \Delta T_{\mathrm{recover},95}
    &=Q_{0.95}\!\left(
      \left\{T_{\mathrm{chain},e}^{\mathrm{C3\text{-}R}}
      -T_{\mathrm{chain},p(e)}^{\mathrm{C2\text{-}A}}:
      e\ \text{触发恢复}\right\}
    \right).
  \end{aligned}
  \label{eq:additional-latency}
\end{equation}
第一项由相同观测、相同\(s=s_{\mathrm{middle}}\)的RTC开关配对回放获得，第二项比较同一请求的RTC单分支与加入固定参考噪声分支及残差计算的批处理，两项在每台机器人上各完成不少于1000对请求；第三项只在触发恢复回合中单列恢复对整链周期的影响。三者不合并为一个“附加时延”。C3-R还直接报告端到端推理时延和截止期违约率，以确认加入监测后仍满足精细阶段期限。该方案仅借鉴FAIL-Detect\cite{xu2025faildetect}以成功数据标定失败分数的思路，不继承其保形阈值保证，也不额外训练分数网络。

\begin{figure}[htbp]
  \centering
  \begin{tikzpicture}[
  font=\small,
  >={Stealth[length=2.2mm]},
  data/.style={draw=teal!65!black,rounded corners=2pt,fill=teal!7,
    minimum height=15mm,text width=30mm,align=center,inner sep=2.5pt},
  process/.style={draw=blue!65!black,rounded corners=2pt,fill=blue!7,
    minimum height=15mm,text width=30mm,align=center,inner sep=2.5pt},
  robot/.style={draw=orange!75!black,rounded corners=2pt,fill=orange!10,
    minimum height=15mm,text width=30mm,align=center,inner sep=2.5pt},
  decide/.style={draw=red!55!black,rounded corners=2pt,fill=red!4,
    minimum height=15mm,text width=30mm,align=center,inner sep=2.5pt},
  recover/.style={draw=violet!65!black,rounded corners=2pt,fill=violet!7,
    minimum height=15mm,text width=30mm,align=center,inner sep=2.5pt},
  calib/.style={draw=gray!65,rounded corners=5pt,fill=gray!8,
    text width=30mm,align=center,inner sep=2.5pt,font=\footnotesize},
  arrow/.style={->,thick,blue!70!black},
  dataarrow/.style={->,thick,teal!70!black},
  alarm/.style={->,thick,red!65!black},
  feedback/.style={->,thick,orange!80!black,rounded corners=3pt},
  influence/.style={->,densely dashed,gray!70,rounded corners=2pt},
  databus/.style={thick,teal!70!black}
]
  \node[data] (prev) {\textbf{前一RTC动作块}\\时间对齐的\\未执行尾部};
  \node[process,right=6mm of prev] (rtc) {\textbf{RTC连续衔接}\\冻结\(d_i\)步前缀\\引导重叠区并补全};
  \node[robot,right=6mm of rtc] (exec) {\textbf{执行与观测}\\执行RTC动作块\\获取新观测};
  \node[data,below=8mm of prev] (raw) {\textbf{当前原始提议}\\新观测+\(\pi_{0.5}\)\\固定参考噪声};
  \node[process,below=8mm of rtc] (inc) {\textbf{跨块一致性残差}\\归一化\(r_i\)\\滑动均值\(\widetilde r_i\)};
  \node[decide,below=8mm of exec] (dec) {\textbf{报警判定}\\\(\widetilde r_i>\eta\)\\连续\(n_{\mathrm{hold}}\)次？};
  \node[recover,below=8mm of dec] (rec) {\textbf{模板恢复}\\停止当前块\\退至\(z_{\mathrm{safe}}\)→回入口位姿};
  \node[calib,below=8mm of raw] (thr) {\textbf{阈值标定（离线）}\\成功回合的全程最大值\\取95分位};
  \node[calib,below=8mm of inc] (shadow) {\textbf{离线检测评价}\\回放记录的C2回合\\只记录报警，不改变动作};
  \node[recover,below=8mm of rec] (retry) {\textbf{重试当前子任务一次}\\成功则继续后序任务\\整链不再允许第二次恢复};
  \node[decide,below=8mm of retry] (stop) {\textbf{安全停机}\\条件不满足或再次报警\\交人工处理};

  % 两份时间对齐的动作提议先在节点间的空白带汇入数据总线，
  % 再分别送往RTC衔接与一致性计算，避免交叉连线穿过文本块。
  \coordinate (bus-top) at ($(prev.east)!0.5!(rtc.west)$);
  \coordinate (bus-bottom) at ($(raw.east)!0.5!(inc.west)$);
  \draw[databus] (prev.east) -- (bus-top) -- (bus-bottom) -- (raw.east);
  \draw[dataarrow] (bus-top) -- (rtc.west);
  \draw[dataarrow] (bus-bottom) -- (inc.west);
  \draw[arrow] (rtc) -- (exec);
  \draw[arrow] (inc) -- (dec);
  % 离线信息只从节点间通道接入，不与在线恢复主链共线。
  \draw[influence] (thr.north) -- ++(0,4mm) -|
    node[pos=0.34,above,font=\footnotesize,fill=white,inner sep=1pt]{既定阈值}
    ($(dec.south west)+(5mm,0)$);
  \draw[influence] (inc.south) -- (shadow.north);
  \draw[arrow] (dec.north) -- node[right,font=\footnotesize]{否：继续} (exec.south);
  \draw[alarm] (dec) -- node[right,font=\footnotesize,fill=white,inner sep=1pt]{C3在线：是} (rec);
  \draw[arrow] (rec) -- (retry);
  \draw[alarm] (retry) -- node[right,font=\footnotesize]{不安全/再次报警} (stop);
  \draw[feedback] (retry.east) -- ++(6mm,0) |-
    node[pos=0.24,right,font=\footnotesize,align=left,fill=white,inner sep=1pt]{恢复后\\重新推理}
    (exec.east);
  \draw[feedback] (exec.north) -- ++(0,6mm) -|
    node[pos=0.24,above,font=\footnotesize,fill=white,inner sep=1pt]{新观测→下一次\(\pi_{0.5}\)推理}
    ([xshift=-7mm]raw.west) -- (raw.west);
  \node[draw=gray!45,densely dashed,rounded corners=3pt,fit=(thr)(shadow),
    inner sep=2.5mm,label={[font=\footnotesize,text=gray!70]below:离线标定与检测评价}] {};
\end{tikzpicture}

  \caption{基于跨动作块一致性残差的失败检测与模板恢复流程}
  \label{fig:ensemble-monitor-recovery}
\end{figure}
\FloatBarrier

\subsection{整体验证方案}

分层鲁棒增强框架按表\ref{tab:stacked-configs}依次验证。C0、C1与C2-A完成平移、航向和自然停靠三套主测试；C2-RS、C2-RM、C2-RL用于固定时域对照，C3-D顺序回放冻结的独立C2-A轨迹进行离线影子报警评价，C3-R在匹配的新回合中启用恢复。五子任务链的整链与条件成功率按式\eqref{eq:task-success-metrics}计算，其长链评价思想参照CALVIN\cite{mees2022calvin}，但完成判据、分母和统计单位均由本文明确定义。所有机器人、配置、残差档位和方向序列使用统一日志时钟，逐次记录配置与动作块编号、阶段\(\phi\)、执行时域、RTC开关、\(r_i\)、\(\widetilde r_i\)、\(g_i\)、\(\eta_b\)、推理请求与返回、截止期违约、导航与参考残差、报警可用、恢复资格判断、停止与恢复起止、\(Y_{e,m}\)、失败原因、人工接管及安全事件，并关联机器人、日期、工位、模型检查点和训练种子。

示范只由一台源机器人采集，按日期分成训练集与超参数验证集；另外三台机器人不进入策略训练。四台机器人的成功阈值标定集均与最终测试日期和回合隔离。主结果分别报告源机器人留出日期和三台未参与训练机器人的结果，所有策略主比较均冻结同一组三种子检查点，不进行按机器人微调。单个二项比例报告Wilson 95\%置信区间；跨配置差值、\(R_{80}\)和归一化曲线面积采用按机器人、训练种子、日期和四停靠方向序列分层的10000次簇自举区间；多个主要比较采用Holm校正。实验顺序按机器人--日期区组随机化，回合是成功率与误报率的统计单位，单个回合内的控制步或推理时刻不作为独立重复。正式试验前依据C0独立预试验的基线率和机器人内相关性，对10个百分点优效、5个百分点非劣界和检测目标分别按双侧\(\alpha=0.05\)、功效0.80计算所需回合数；计算值高于各节明示下限时扩充至计算值，并把预试验数据、功效脚本和最终样本数一并冻结。

\begin{longtable}{@{}p{1.4cm}p{3.6cm}p{8.4cm}@{}}
  \caption{整体验证的叠加配置}\label{tab:stacked-configs}\\
  \toprule
  配置 & 组成 & 主要考察指标 \\
  \midrule
  \endfirsthead
  \multicolumn{3}{c}{续表~\thetable}\\
  \toprule
  配置 & 组成 & 主要考察指标 \\
  \midrule
  \endhead
  C0 & \(\pi_{0.5}\)固定执行时域功能基线 & 式\eqref{eq:task-success-metrics}、残差曲线、\(R_{80}\)与\(\operatorname{AUC}_{\mathrm{rob}}\) \\
  C1 & C0加入随机停靠示范和停靠残差估计状态 & 与C0及消融(b)比较\(R_{80}\)、归一化曲线面积和成功率差 \\
  C2-R & C1分别加入短、中、长固定时域及相同RTC（RS/RM/RL） & 固定时域下的条件成功率、推理频率、时延和截止期违约率 \\
  C2-A & C1加入阶段自适应时域及RTC & 与短、长固定RTC组比较条件成功率、容忍边界和GPU执行时间 \\
  C3-D & 离线回放C2-A轨迹并运行冻结检测器 & 回合召回率、成功回合误报率、每小时非窗口报警与提前量分布 \\
  C3-R & C2-A加入冻结检测器和一次模板恢复 & 首次及最终成功率、条件恢复成功率、人工接管、安全事件与恢复周期 \\
  \bottomrule
\end{longtable}

\subsection{可行性分析}

本课题的可行性建立在已完成接口和方法的可分步实现上。其一，第\ref{sec:completed}节所述VR示范采集、OpenPI数据组织、\(\pi_{0.5}\)任务微调与动作块执行接口已经贯通，可作为功能基线；五子任务整链成功率及残差敏感性仍按本节协议重新统计。其二，三项方法不改变策略网络主体结构：残差输入扩展状态接口，OpenPI任务微调已可在课题分配的一张A100 80G显卡上运行；分阶段执行、RTC和监测位于推理侧；检测阈值不使用失败训练数据。其三，导航目标可实施程序化注入，激光雷达定位提供策略输入，新增的独立工位位姿测量负责实际残差归档，三者职责分离。其四，四台机器人按区组并行执行时，每个主曲线配置的三个种子合计至少包含528条轴向残差链和96条自然停靠链；每个固定时域配置至少240条链，C3-R至少240对匹配链，C3-D优先复用与标定集隔离的C2-A轨迹并在失败数不足时按冻结压力档位补足。C0预试验后用实测单链时长核算总机器人小时；若功效扩充后的总量超过四台机器人在进度表内可提供的有效工时，则在正式测试前延长实验时段，不削减三种子主比较或事后放宽统计判据。可能的风险与应对措施见第\ref{sec:risks}节。

\section{预期达到的目标}

本课题预期在已贯通的OpenPI \(\pi_{0.5}\)接口链路基础上，建立基于工位停靠残差传递链的分层鲁棒增强框架，并在四台轮式双臂机器人和真实的PCBA四工位五子任务链上完成验证。产线周期上限\(T_{\mathrm{takt}}\)采用现场有效版本工艺文件中从首次工位导航请求到成品放置稳定的秒级上限；若现行文件尚无该数值，则由产线工艺负责人依据上下游供料周期在2026年10月前给出书面上限及版本号，未获得数值前不启动周期验收。1厘米与\(1^\circ\)对应一档可分辨残差改善，10个百分点定义为成功率最小实际重要差异，5个百分点定义为非劣界，20\%定义为值得引入自适应调度的最小GPU时间节省；这些口径与控制保护阈值共同写入冻结协议，并经第\ref{subsec:task-definition}节的预试验和功效计算核查。若现场工艺规范更严格，则采用较严格值且不在查看测试结果后调整。

对于条件化学习方法，C1相对C0的平移\(R_{80}\)目标增加不少于1厘米、航向\(R_{80}\)目标增加不少于\(1^\circ\)，两条曲线的\(\operatorname{AUC}_{\mathrm{rob}}\)差值95\%置信下界均大于0；在3厘米和\(3^\circ\)档位，C1整链成功率的分层簇自举95\%下界目标不低于0.80。显式残差配置(c)相对相同随机数据的配置(b)，在上述两个训练边界档位的成功率目标分别提高不少于10个百分点，且分层簇自举区间下界大于0。各配置从同一预训练检查点按相同预算独立微调三个种子，而不表述为同一个模型。

对于自适应控制方法，五个代表档位等权汇总时，C2-A相对全程短时域C2-RS的精细子任务条件成功率采用5个百分点的非劣界，即差值95\%区间下界不低于\(-0.05\)，同时每链GPU执行时间\(C_{\mathrm{GPU}}\)至少降低20\%；相对全程长时域C2-RL，精细子任务条件成功率目标提高不少于10个百分点且区间下界大于0。C2-RM只作次要趋势分析。C2-A的平移或航向\(R_{80}\)至少有一项相对C1增加一个档位，并满足\(Q_{0.95}(\tau_{\mathrm{inf}})\le s_{\mathrm{fine}}\Delta t\)及截止期违约率不高于1\%。

对于运行时容错方法，C3-D的回合召回率点估计目标不低于90\%、Wilson 95\%下界不低于80\%；成功回合误报率点估计目标不高于5\%、Wilson 95\%上界不高于10\%。及时检出的提前量要求\(Q_{0.05}(T_{\mathrm{lead}})>T_{\mathrm{stop,safe}}\)。C3-R的条件恢复成功率点估计目标不低于50\%、Wilson 95\%下界不低于35\%，最终整链成功率相对匹配C2-A提高不少于10个百分点且分层簇自举差值区间下界大于0。监测开销要求\(\Delta\tau_{\mathrm{mon},95}\le0.1\Delta t\)，并直接要求\(Q_{0.95}(\tau_{\mathrm{C3\text{-}R}})\le s_{\mathrm{fine}}\Delta t\)、C3-R截止期违约率不高于1\%；未触发恢复回合相对匹配C2-A的整链周期增幅单侧95\%上界不超过5\%，触发恢复回合的周期95分位数不超过\(T_{\mathrm{takt}}\)。恢复诱发的不可逆物料或夹具损伤目标为0；否则即使成功率达标，也不判定运行时容错方案通过验收。

预期成果形式包括：面向OpenPI \(\pi_{0.5}\)的工位停靠残差鲁棒性改进及其真机验证；主配置与桥接配置在统一残差协议下的真机数据、容忍曲线、检测与恢复结果；包含冻结参数、试验次数、置信区间、失败类型和适用边界的评价报告；以及以上述内容为主体的硕士学位论文。

\section{已完成的学术研究/实践工作}
\label{sec:completed}

\subsection{实验平台与坐标系定义}

本研究已在轮式双臂机器人平台上完成示范采集与策略执行所需的基础接口建设。机器人上半身由左右两条七自由度机械臂、四自由度腰部和二自由度头部组成，共包含20个转动关节，并配置末端夹爪、头部相机和腕部相机。现阶段的双臂操作以左右机械臂末端笛卡尔位姿和夹爪状态作为核心状态与动作变量，执行模型将腰部、头部和移动底盘位姿作为固定参数，仅更新左右机械臂的关节状态。移动底盘由激光雷达定位与传统导航算法在工位间移动并停靠，停靠后锁定，操作阶段仅由双臂执行。

为统一描述机器人、工位与虚拟现实（Virtual Reality, VR）设备之间的空间关系，本文以 \(\alpha\in\{L,R\}\) 区分左、右机械臂，并定义以下坐标系：世界坐标系记为 \(\{O\}\)，机器人基座坐标系记为 \(\{B\}\)，左右末端执行器坐标系记为 \(\{E_\alpha\}\)，第 \(k\) 个工位固连于环境且与示范时名义基座位姿重合的停靠坐标系记为 \(\{S_k\}\)；VR参考坐标系记为 \(\{V\}\)，左右手控制器坐标系记为 \(\{H_\alpha\}\)。其中，\({}^{A}\mathbf{T}_{C}\) 表示坐标系 \(\{C\}\) 相对于坐标系 \(\{A\}\) 的位姿，即把 \(\{C\}\) 中的齐次坐标变换到 \(\{A\}\) 中。该表示写为
\begin{equation}
  {}^{A}\mathbf{T}_{C}
  =
  \begin{bmatrix}
    {}^{A}\mathbf{R}_{C} & {}^{A}\mathbf{p}_{C} \\
    \mathbf{0}_{1\times 3} & 1
  \end{bmatrix},
  \label{eq:homogeneous-transform}
\end{equation}
式中，\({}^{A}\mathbf{R}_{C}\in\mathrm{SO}(3)\) 为旋转矩阵，\({}^{A}\mathbf{p}_{C}\in\mathbb{R}^{3}\) 为坐标系原点的位置向量。若空间点在 \(\{C\}\) 中的齐次坐标为 \(\bar{\mathbf p}^{C}\)，则其在 \(\{A\}\) 中的表示为 \(\bar{\mathbf p}^{A}={}^{A}\mathbf T_C\bar{\mathbf p}^{C}\)。连续坐标变换满足
\begin{equation}
  {}^{A}\mathbf{T}_{D}
  =
  {}^{A}\mathbf{T}_{C}\,{}^{C}\mathbf{T}_{D},
  \qquad
  \left({}^{A}\mathbf{T}_{C}\right)^{-1}
  =
  \begin{bmatrix}
    \left({}^{A}\mathbf{R}_{C}\right)^{\mathrm T}
    &
    -\left({}^{A}\mathbf{R}_{C}\right)^{\mathrm T}{}^{A}\mathbf{p}_{C} \\
    \mathbf{0}_{1\times 3} & 1
  \end{bmatrix}.
  \label{eq:transform-compose-inverse}
\end{equation}
式\eqref{eq:homogeneous-transform}和式\eqref{eq:transform-compose-inverse}构成VR增量映射、工位坐标系与前文停靠残差定义的统一表示基础。基座系与末端系之间的正逆运动学由平台现有笛卡尔控制模块提供，本文不再展开。

\subsection{基于VR遥操作的示范数据采集}

在上述坐标关系基础上，本研究已构建基于VR设备的双臂遥操作与示范采集链路。系统读取头显及左右手控制器的空间位姿和交互信号，在遥操作回合开始时分别记录控制器初始位姿
\({}^{V}\mathbf T_{H_\alpha}(t_0)\) 与机器人末端初始位姿
\({}^{B}\mathbf T_{E_\alpha}(t_0)\)。固定标定旋转
\({}^{B}\mathbf R_V\) 用于统一VR参考系与机器人基座系的轴向约定，比例系数 \(\lambda_\alpha\) 用于调节操作者手部平移与机器人末端平移的尺度。由此，控制器相对初始位置的增量被映射为
\begin{equation}
  {}^{B}\mathbf p_{E_{\alpha,d}}(t)
  =
  {}^{B}\mathbf p_{E_\alpha}(t_0)
  +
  \lambda_\alpha{}^{B}\mathbf R_V
  \left(
    {}^{V}\mathbf p_{H_\alpha}(t)
    -
    {}^{V}\mathbf p_{H_\alpha}(t_0)
  \right).
  \label{eq:vr-position-mapping}
\end{equation}
该增量形式以遥操作开始时的实际末端位置为基准，VR空间的绝对坐标用于计算控制器相对初始位姿的运动增量。

控制器姿态首先相对于初始姿态计算，再经标定旋转映射到机器人基座系。末端目标姿态表示为
\begin{equation}
  \begin{aligned}
    {}^{B}\mathbf R_{E_{\alpha,d}}(t)
    ={}^{B}\mathbf R_V
    \left[
      {}^{V}\mathbf R_{H_\alpha}(t)
      \left({}^{V}\mathbf R_{H_\alpha}(t_0)\right)^{\mathrm T}
    \right]
    {}^{V}\mathbf R_B
    {}^{B}\mathbf R_{E_\alpha}(t_0),
  \end{aligned}
  \label{eq:vr-orientation-mapping}
\end{equation}
其中 \({}^{V}\mathbf R_B=\left({}^{B}\mathbf R_V\right)^{\mathrm T}\)。式\eqref{eq:vr-position-mapping}与式\eqref{eq:vr-orientation-mapping}分别生成末端目标位置和姿态，两者组成
\({}^{B}\mathbf T_{E_{\alpha,d}}(t)\)，再由已有笛卡尔控制接口执行。手柄交互量 \(u_\alpha(t)\) 通过既定映射
\begin{equation}
  g_{\alpha,d}(t)=\mathcal G_\alpha\!\left(u_\alpha(t)\right)
  \label{eq:gripper-mapping}
\end{equation}
转换为左右夹爪的开合指令。映射函数 \(\mathcal G_\alpha(\cdot)\) 的触发阈值与夹爪行程参数由平台配置确定。

机器人示范数据通常以一次完整的任务回合作为基本组织单元，将视觉观测、机器人状态与动作按时间步对齐，并在回合层关联任务指令；统一的数据语义有助于后续按不同观测窗口和动作空间构造训练样本\cite{openx2024rtx}。参照这一组织方式，本研究将当前采集的数据划分为任务、回合和时间步三个层级：任务层以自然语言描述界定操作目标，回合层对应一次完整的遥操作过程，时间步层则以动作时刻为索引，关联头部与左右腕部图像、双臂末端状态、夹爪状态及相应动作。第 \(n\) 个任务回合可表示为
\begin{equation}
  \tau^{(n)}
  =
  \left\{
    \left(
      t_k,
      I^{\mathrm{head}}_k,
      I^{\mathrm{wrist},L}_k,
      I^{\mathrm{wrist},R}_k,
      \mathbf x_k,
      \mathbf a_k
    \right)
  \right\}_{k=1}^{K_n},
  \label{eq:demonstration-trajectory}
\end{equation}
其中，\(\mathbf x_k\) 包含左右末端位姿与夹爪状态，\(\mathbf a_k\) 包含对应时刻的双臂笛卡尔动作与夹爪指令。任务语言描述作为策略的语义条件与相应示范回合关联；一个示范回合对应一次低层子任务执行，正式评价中的五子任务链则由状态机串联这些子任务。上述层级将任务语义与单次执行序列分开，同时保持观测、状态和动作在回合内部逐时对应，可直接为后续视觉--语言--动作策略构造单步样本或动作块样本。现有示范均在名义停靠位姿下采集，第\ref{subsec:improvement-one}节将把采集扩展为随机停靠，并在回合层同步记录所属停靠事件的\(\boldsymbol\delta_k^{\mathrm{cmd}}\)、\(\widehat{\boldsymbol\delta}_k^{\mathrm{loc}}\)与\(\boldsymbol\delta_k^{\mathrm{ref}}\)。

以“将印制电路板组件（Printed Circuit Board Assembly，PCBA）放置到百叶车指定承载位置”为例，该回合以右臂为操作臂，左臂保持固定。机器人在夹爪夹持状态下携带PCBA接近百叶车，调整末端位置和姿态完成放置，随后松开夹爪并撤离放置区域。该回合共包含177个时间对齐样本；每个时间步包含三路分辨率为\(240\times320\)的彩色图像，以及14维双臂末端状态和14维双臂动作。状态与动作均按左右臂各7维组织，分别由三维位置、滚转--俯仰--偏航姿态和夹爪状态或指令组成。该样例表明，现有采集链路能够将多视角观测、机器人本体状态与遥操作动作组织为式\eqref{eq:demonstration-trajectory}所定义的完整回合。

图\ref{fig:pcba-louver-cart-demo}给出了该回合的起始、中间和结束观测，以及右手末端的三维运动轨迹。上部三幅图分别对应第0、88和176帧：前两帧处于夹持状态，结束帧处于松开状态；下部曲线以颜色和线型区分夹爪语义状态，以箭头表示工具坐标系\(+Z\)轴方向。轨迹在第120帧由夹持切换为松开，使视觉观测、末端位姿和PCBA释放事件能够在同一时间轴上对应。

\begin{figure}[htbp]
  \centering
  \includegraphics[width=\textwidth]{right_hand_visualization/right_hand_overview.png}
  \caption{将PCBA放置到百叶车任务的示例观测、右手轨迹与夹爪状态}
  \label{fig:pcba-louver-cart-demo}
\end{figure}

图\ref{fig:preliminary-system-loop}概括了已完成系统中数据采集与策略执行之间的接口关系。VR位姿经过增量映射后进入笛卡尔控制模块，机器人动作及多视角观测由同步记录模块组织为示范轨迹；示范数据经OpenPI数据接口适配后用于\(\pi_{0.5}\)微调，部署阶段由任务状态机提供动态低层语言指令，模型接收新观测并通过同一笛卡尔接口输出动作块。由此，“VR数采--\(\pi_{0.5}\)微调--OpenPI推理--真机执行”的接口链路已经贯通并完成单回合功能验证，示范阶段与推理阶段共享机器人状态、动作定义和底层执行链路；该结果不代表五子任务整链性能。

\begin{figure}[htbp]
  \centering
  \begin{tikzpicture}[
  font=\small,
  >={Stealth[length=2.2mm]},
  source/.style={draw=violet!65!black,rounded corners=2pt,fill=violet!7,
    minimum height=14mm,text width=25mm,align=center,inner sep=2.5pt},
  process/.style={draw=blue!65!black,rounded corners=2pt,fill=blue!7,
    minimum height=14mm,text width=27mm,align=center,inner sep=2.5pt},
  robot/.style={draw=orange!75!black,rounded corners=2pt,fill=orange!10,
    minimum height=14mm,text width=28mm,align=center,inner sep=2.5pt},
  data/.style={draw=teal!65!black,rounded corners=2pt,fill=teal!7,
    minimum height=14mm,text width=30mm,align=center,inner sep=2.5pt},
  arrow/.style={->,thick,blue!70!black},
  dataarrow/.style={->,thick,teal!70!black},
  trainarrow/.style={->,thick,densely dashed,violet!70!black},
  feedback/.style={->,thick,orange!80!black,rounded corners=3pt}
]
  \node[source] (vr) {\textbf{VR遥操作}\\双手位姿\\与夹爪输入};
  \node[process,right=5mm of vr] (mapping) {\textbf{坐标映射}\\参考位姿标定\\与运动增量};
  \node[process,right=5mm of mapping] (controller) {\textbf{笛卡尔接口}\\逆运动学\\与轨迹约束};
  \node[robot,right=5mm of controller] (robot) {\textbf{双臂机器人}\\末端动作执行\\与多视角感知};

  \draw[arrow] (vr) -- (mapping);
  \draw[arrow] (mapping) -- (controller);
  \draw[arrow] (controller) -- (robot);

  \node[data,below=13mm of robot] (record) {\textbf{同步记录}\\图像·末端状态\\动作·时间戳};
  \node[data,left=7mm of record] (dataset) {\textbf{示范轨迹}\\按任务回合组织\\语言条件关联};
  \node[source,left=7mm of dataset] (policy) {\textbf{OpenPI \(\pi_{0.5}\)}\\离线训练\\与模型部署};
  \node[process,left=7mm of policy] (online) {\textbf{在线推理}\\新观测+动态指令\\输出动作块};

  \draw[dataarrow] (robot) -- (record);
  \draw[dataarrow] (record) -- (dataset);
  \draw[trainarrow] (dataset) -- (policy);
  \draw[arrow] (policy) -- (online);
  \draw[arrow] (online.north) -- ++(0,6mm) -| (controller.south);
  \draw[feedback] (robot.east) -- ++(5mm,0) |-
    ([yshift=-7mm]online.south) -- (online.south);
\end{tikzpicture}

  \caption{已贯通的VR数采--\(\pi_{0.5}\)微调--OpenPI推理--真机执行接口链路}
  \label{fig:preliminary-system-loop}
\end{figure}
\FloatBarrier

\subsection{基于OpenPI的\texorpdfstring{\(\pi_{0.5}\)}{pi0.5}微调与真机闭环}

在示范采集链路基础上，现有系统已完成OpenPI数据适配、\(\pi_{0.5}\)任务示范微调、在线推理服务接入和真机动作块执行。当前部署使用OpenPI公开的\(\pi_{0.5}\)模型及流匹配动作头：模型接收头部与左右腕部图像、双臂本体状态和任务状态机提供的低层语言指令，生成连续动作块，再由笛卡尔控制接口执行并在动作块边界回传新观测。这里仅声称数据采集--微调--动作块执行的软件与真机接口链路已经贯通；五子任务整链成功率、工位停靠残差敏感性和相对性能尚待C0--C3实验统计。

为说明当前实现所用模型的算法基础，原始\(\pi_{0.5}\)在\(\pi_0\)连续动作生成框架的基础上引入异构协同训练与语义子任务预测，将长时程任务中的高层语义决策与低层动作块生成组织在统一模型中，其结构及训练流程如图\ref{fig:pi05-model-overview}所示\cite{black2025pi05}。

\begin{figure}[htbp]
  \centering
  \includegraphics[width=0.98\textwidth]{pi05-model-overview.png}
  \caption{\(\pi_{0.5}\)模型结构及两阶段训练流程\cite{black2025pi05}}
  \label{fig:pi05-model-overview}
\end{figure}

图\ref{fig:pi05-model-overview}左侧为预训练阶段。模型以SigLIP视觉编码器和Gemma语言模型组成的预训练VLM为主干，将多模态网络数据、异构机器人数据和任务提示统一表示为序列，并同时学习语义子任务、离散动作、图像描述及目标框等输出。机器人动作在该阶段由FAST分词器压缩为离散词元，使动作预测能够与语言建模任务共同采用自回归的下一词元预测目标。图中右侧为后训练与推理阶段：预训练得到的VLA模型继续承担图像与语言理解，并在主干之外加入参数量较小的动作专家；动作专家接收连续动作的含噪中间状态，采用条件流匹配逐步生成连续动作块。由此，模型既保留预训练阶段获得的视觉与语义能力，又能以非自回归方式输出适用于实时控制的连续动作。

在层级推理过程中，VLA主干首先根据当前多视角观测和高层任务指令\(g_t\)自回归生成语义子任务\(\widehat{\ell}_t\)，例如将“整理卧室”转化为“拿起枕头”；随后，该子任务作为低层语言条件输入动作专家，后者由随机噪声生成长度为\(H\)的连续动作块。原模型的联合分布可分解为
\begin{equation}
  \pi_{\boldsymbol\theta}
  \left(\mathbf A_t,\widehat{\ell}_t\mid\mathbf z_t,g_t\right)
  =
  \pi_{\boldsymbol\theta}
  \left(\mathbf A_t\mid\mathbf z_t,\widehat{\ell}_t\right)
  \pi_{\boldsymbol\theta}
  \left(\widehat{\ell}_t\mid\mathbf z_t,g_t\right),
  \label{eq:pi05-hierarchical-factorization}
\end{equation}
其中，\(\mathbf z_t\)表示图像与机器人本体状态，\(g_t\)表示原模型的高层任务指令。式\eqref{eq:pi05-hierarchical-factorization}的第二项对应原模型的子任务预测，第一项对应动作专家的连续动作生成。上述异构预训练与高层子任务预测属于原始\(\pi_{0.5}\)的模型能力；本课题不宣称复现其预训练过程，也不把高层子任务预测作为研究变量。当前真机接口只使用OpenPI公开模型及流匹配动作头，由外部任务状态机根据五个子任务直接提供动态低层语言指令\(\ell_t\)。为与现有推理接口保持一致，时刻 \(t\) 的策略输入写为
\begin{equation}
  \mathbf o_t
  =
  \left\{
    I^{\mathrm{head}}_t,
    I^{\mathrm{wrist},L}_t,
    I^{\mathrm{wrist},R}_t,
    \mathbf x_t,
    \ell_t
  \right\},
  \label{eq:openpi-observation}
\end{equation}
其中，\(\ell_t\) 为运行时给定的动态任务语言指令。策略一次推理输出长度为 \(H\) 的动作块：
\begin{equation}
  \widehat{\mathbf A}_t
  =
  \pi_{\boldsymbol\theta}(\mathbf o_t)
  =
  \left[
    \widehat{\mathbf a}_{t},
    \widehat{\mathbf a}_{t+1},
    \ldots,
    \widehat{\mathbf a}_{t+H-1}
  \right].
  \label{eq:openpi-action-chunk}
\end{equation}
动作块中的每个元素均按照前述双臂笛卡尔动作与夹爪指令接口解释。机器人执行其中时域为\(s\)的近期动作，且\(1\leq s\leq H\)；完成该段动作后，系统重新获取多视角图像和末端状态，并结合当前语言指令构造下一轮输入。该过程在动作块边界形成“状态感知--策略推理--动作块执行--状态更新”的视觉闭环，而每次提交的\(s\)个控制步在策略层面不接收新观测。现有部署中\(s\)为全任务统一的固定值，第\ref{subsec:improvement-two}节将其改为按阶段取值并加入RTC衔接。

\paragraph{面向\(\pi_{0.5}\)的模仿学习推导。}将式\eqref{eq:demonstration-trajectory}所示的专家示范按照长度为 \(H\) 的滑动窗口组织，可得动作块示教数据集
\begin{equation}
  \mathcal D_{\mathrm E}
  =
  \left\{
    \left(\mathbf o_t^{(n)},\mathbf A_t^{\mathrm E,(n)}\right)
  \right\}_{n,t},
  \qquad
  \mathbf A_t^{\mathrm E}
  =
  \left[
    \mathbf a_t^{\mathrm E},\ldots,
    \mathbf a_{t+H-1}^{\mathrm E}
  \right].
  \label{eq:pi-demonstration-dataset}
\end{equation}
行为克隆的基本目标是在专家数据分布上最大化条件动作块的对数似然，即
\begin{equation}
  \boldsymbol\theta_{\mathrm{IL}}
  =
  \arg\min_{\boldsymbol\theta}\mathcal L_{\mathrm{IL}}(\boldsymbol\theta),
  \qquad
  \mathcal L_{\mathrm{IL}}(\boldsymbol\theta)
  =
  -\mathbb E_{(\mathbf o_t,\mathbf A_t^{\mathrm E})\sim\mathcal D_{\mathrm E}}
  \left[
    \log \pi_{\boldsymbol\theta}
    \left(\mathbf A_t^{\mathrm E}\mid\mathbf o_t\right)
  \right].
  \label{eq:pi-il-nll}
\end{equation}
该目标使策略分布逼近专家在给定多视角图像、机器人状态和语言指令条件下的动作分布；动作块联合建模还保留了同一操作片段内各时刻动作之间的相关性。

\(\pi_{0.5}\)采用条件流匹配表示连续动作块分布\cite{black2024pi0,black2025pi05}。对归一化后的专家动作块 \(\mathbf A_t^{\mathrm E}\)，采样高斯噪声 \(\boldsymbol\epsilon\sim\mathcal N(\mathbf 0,\mathbf I)\) 和流时间 \(\tau\in[0,1]\)，构造从噪声到数据的线性概率路径
\begin{equation}
  \mathbf A_t^{\tau}
  =
  (1-\tau)\boldsymbol\epsilon
  +
  \tau\mathbf A_t^{\mathrm E},
  \qquad
  \mathbf u_t^{\tau}
  =
  \frac{\partial\mathbf A_t^{\tau}}{\partial\tau}
  =
  \mathbf A_t^{\mathrm E}-\boldsymbol\epsilon.
  \label{eq:pi-flow-path}
\end{equation}
设动作专家输出的条件向量场为
\(\mathbf v_{\boldsymbol\theta}(\mathbf A_t^{\tau},\mathbf o_t,\tau)\)，则模仿微调可写为条件流匹配损失
\begin{equation}
  \mathcal L_{\mathrm{FM}}(\boldsymbol\theta)
  =
  \mathbb E_{\substack{
    (\mathbf o_t,\mathbf A_t^{\mathrm E})\sim\mathcal D_{\mathrm E},\;
    \boldsymbol\epsilon\sim\mathcal N(\mathbf 0,\mathbf I),\;
    \tau\sim p(\tau)}}
  \left[
    \left\|
      \mathbf v_{\boldsymbol\theta}
      (\mathbf A_t^{\tau},\mathbf o_t,\tau)
      -
      (\mathbf A_t^{\mathrm E}-\boldsymbol\epsilon)
    \right\|_2^2
  \right].
  \label{eq:pi-flow-matching-loss}
\end{equation}
最小化式\eqref{eq:pi-flow-matching-loss}使网络学习式\eqref{eq:pi-flow-path}的条件速度场。推理时从
\(\mathbf A_t^{0}\sim\mathcal N(\mathbf 0,\mathbf I)\) 出发，沿
\begin{equation}
  \frac{\mathrm d\mathbf A_t^{\tau}}{\mathrm d\tau}
  =
  \mathbf v_{\boldsymbol\theta}
  (\mathbf A_t^{\tau},\mathbf o_t,\tau)
  \label{eq:pi-flow-ode}
\end{equation}
由 \(\tau=0\) 正向积分至 \(\tau=1\)，即可得到条件动作块样本
\(\widehat{\mathbf A}_t=\mathbf A_t^1\)。以步长 \(\Delta\tau>0\) 进行欧拉离散时，单步更新为
\begin{equation}
  \mathbf A_t^{\tau+\Delta\tau}
  =
  \mathbf A_t^{\tau}
  +
  \Delta\tau\,
  \mathbf v_{\boldsymbol\theta}
  (\mathbf A_t^{\tau},\mathbf o_t,\tau).
  \label{eq:pi-flow-euler}
\end{equation}
因此，式\eqref{eq:pi-flow-matching-loss}给出了示教动作对模型参数的监督关系，式\eqref{eq:pi-flow-ode}和式\eqref{eq:pi-flow-euler}则给出了从随机噪声生成连续动作块的推理过程。

\subsection{阶段性工作小结}

本节明确了已完成工作与待开展研究的边界。已完成部分包括机器人平台与坐标接口、VR遥操作和示范数据采集、OpenPI数据适配、\(\pi_{0.5}\)任务示范微调以及“状态感知--OpenPI推理--动作块执行--状态更新”的单回合接口功能验证；激光雷达定位与导航停靠作为平台既有能力可直接调用。尚未开展的部分对应本课题的三项方法：工位停靠残差尚未进入\(\pi_{0.5}\)状态，示范仍在名义停靠下采集，执行时域仍为全任务固定值，推理侧尚未加入RTC连续衔接、跨动作块一致性检测和模板恢复。因尚无统一残差注入协议下的统计数据，本节不对现有接口链路声明整链成功率或性能优势。

\Needspace{12\baselineskip}
\section{进度安排}

研究进度将从2026年7月进入课题时开始计算，至2027年12月完成论文与答辩准备。各阶段按照“时间段--主要任务--阶段产出”编排，并在后期预留补充实验、设备维护和论文修改时间。

\begin{longtable}{@{}p{3.0cm}p{7.0cm}p{3.4cm}@{}}
  \caption{研究进度安排}\label{tab:research-schedule}\\
  \toprule
  时间 & 主要任务 & 阶段产出 \\
  \midrule
  \endfirsthead
  \multicolumn{3}{c}{续表~\thetable}\\
  \toprule
  时间 & 主要任务 & 阶段产出 \\
  \midrule
  \endhead
  2026.07--2026.08 &
  梳理文献与现有平台接口，确定四工位五子任务链及完成、失败、超时与安全判据；制定随机停靠示范采集规范与残差注入协议，完成开题。 &
  开题报告、任务规范与评价协议初稿。 \\
  2026.09--2026.10 &
  实现导航指令、定位估计与独立参考残差的同步记录及工位标定接口，将残差增广状态接入OpenPI \(\pi_{0.5}\)数据格式；完成外部测量不确定度与重复性验收、冻结档位接纳规则，并按规范采集随机停靠示范。 &
  三类残差接口、测量校准与验收记录、版本化示范数据及数据说明。 \\
  2026.11--2026.12 &
  由同一预训练检查点按三个随机种子完成条件化学习三组消融；按每个主配置至少528条轴向链和96条自然链的下限，执行平移0至5厘米、航向0至\(5^\circ\)和自然停靠测试。 &
  条件化学习实验数据、\(R_{80}\)、归一化曲线面积及带置信区间的消融结论。 \\
  2027.01--2027.03 &
  实现阶段滞回、阶段自适应动作块控制与RTC连续衔接；在五个冻结代表档位完成每个固定时域配置至少240条链，并完成C2-A全曲线，记录条件成功率、推理频率、GPU执行时间及截止期违约。 &
  分阶段执行与RTC模块、固定时域桥接对照及运行时统计。 \\
  2027.04--2027.06 &
  在OpenPI推理侧实现跨动作块一致性残差和模板恢复；用独立成功集冻结全回合阈值，先以冻结C2-A轨迹完成C3-D离线影子检测，再完成不少于240对匹配链的C3-R主动恢复及在线截止期验证。 &
  检测与恢复模块、回合召回/误报、提前量、最终成功率、安全事件与分项时延数据。 \\
  2027.07--2027.09 &
  在四台机器人上完成主配置与桥接配置的功效扩充数，补足至少120个成功回合和80个失败回合的检测测试集，整理失败案例并分开报告源机器人和三台目标机器人。 &
  整体验证数据、失败分类与方法适用边界。 \\
  2027.10--2027.12 &
  核对数据、代码和图表，完成学位论文撰写、修改、送审与答辩准备。 &
  学位论文定稿与答辩材料。 \\
  \bottomrule
\end{longtable}

\section{为完成课题所需的条件}

本课题所需条件分为硬件平台、算力、数据、软件与评价环境五类，其中现有接口可直接复用；独立位姿测量、安全设施和评价协议在前两个阶段补齐，RTC、监测与恢复模块按表\ref{tab:research-schedule}的后续阶段实现，如表\ref{tab:conditions}所示。

\begin{longtable}{@{}p{1.8cm}p{5.6cm}p{6.0cm}@{}}
  \caption{完成课题所需的条件}\label{tab:conditions}\\
  \toprule
  类别 & 已具备 & 尚需补充 \\
  \midrule
  \endfirsthead
  \multicolumn{3}{c}{续表~\thetable}\\
  \toprule
  类别 & 已具备 & 尚需补充 \\
  \midrule
  \endhead
  硬件平台 & 四台轮式双臂机器人，各含两条七自由度机械臂、四自由度腰部、二自由度头部、末端夹爪、头部与腕部相机，以及激光雷达定位与导航模块；一台用于示范采集，四台用于实验。 & 工位周边的安全围栏与急停联动；经校准且扩展不确定度不大于2毫米和\(0.2^\circ\)、重复标准差不大于1毫米和\(0.1^\circ\)的独立工位位姿测量或等效装置。 \\
  算力 & 8卡A100 80G服务器，课题分配一张显卡，满足OpenPI \(\pi_{0.5}\)任务微调需求。 & 无。 \\
  数据 & 名义停靠下的VR遥操作示范回合，含PCBA放置样例；数据按任务、回合、时间步三级组织。 & 每子任务不少于50个随机停靠有效示范；相互隔离的训练、验证、成功阈值标定与最终测试数据；双人盲标的失败时刻与安全事件。 \\
  软件 & VR遥操作与同步记录链路；OpenPI \(\pi_{0.5}\)微调与真机部署流水线；导航栈与笛卡尔控制接口。 & 工位停靠残差读取与工位标定接口；残差增广状态的数据格式与归一化；推理侧的阶段判定、分阶段执行、RTC连续衔接与跨块一致性监测模块；模板恢复动作。 \\
  评价环境 & 真实的PCBA四工位五子任务场景：上料百叶车、扫码工位、检测台、下料百叶车。 & 程序化残差注入与参考残差归档工具；覆盖推理、报警、停止、恢复、人工接管及安全事件的统一时钟日志；四台机器人的随机化评价脚本。 \\
  \bottomrule
\end{longtable}

\section{研究过程中可能遇到的问题以及解决的措施}
\label{sec:risks}

三项方法分别依赖残差估计、阶段判定与不一致度信号的可靠性，其实验又受真机资源与本体差异约束。表\ref{tab:risks}按风险来源列出可能遇到的问题、影响以及预防措施与降级方案。

\begin{longtable}{@{}p{2.2cm}p{3.2cm}p{4.0cm}p{4.0cm}@{}}
  \caption{可能遇到的问题与解决措施}\label{tab:risks}\\
  \toprule
  问题 & 影响 & 预防措施 & 降级方案 \\
  \midrule
  \endfirsthead
  \multicolumn{4}{c}{续表~\thetable}\\
  \toprule
  问题 & 影响 & 预防措施 & 降级方案 \\
  \midrule
  \endhead
  激光雷达停靠残差估计存在噪声 & 残差输入本身带误差，削弱条件化学习方法的收益 & 停靠后静止采样多帧取均值；以经校准的独立工位测量量化估计误差 & 若测量不确定度或重复性未达冻结指标，暂停相应轴试验，不以导航指令替代真值 \\
  随机示范未覆盖真实残差分布 & 自然停靠测试落在训练覆盖之外 & 先测自然停靠分布；训练对\(x,y\)各覆盖\(\pm3\)厘米并对航向覆盖\(\pm3^\circ\) & 在不查看最终测试结果的追加采集阶段扩展范围，并重新冻结全部检查点 \\
  阶段判定错误 & 精细阶段误判为自由阶段时开环区间过长，反之推理开销上升 & 任务标签与10/12厘米滞回规则双重判定；预试验后冻结 & 主试验改用全程短RTC配置，并将自适应配置记为未通过而非临时调参 \\
  RTC与监测增加推理时延 & 精细阶段执行时域短于推理延迟 & 异步冻结已确定动作；原始分支批处理；一致性残差向量化；尾部耗尽时受控停止 & 若\(d_\star<\lfloor H/2\rfloor<H-d_\star\)不成立，先优化推理链路并重新预注册时域，不在正式测试中放宽约束 \\
  阈值失准导致误报或漏报 & 误报触发不必要恢复，漏报失去预警意义 & 用成功标定回合的回合最大统计量控制误报预算；参数先于失败测试冻结 & 保持冻结阈值并报告未达标，不用失败标签重新选择工作点 \\
  模板恢复动作引发碰撞 & 回退过程中夹爪或PCBA与槽位干涉 & 回退方向沿工位系竖直方向，路径经离线校核；恢复前先停止运动 & 报警后直接停机交人工，不执行自动恢复 \\
  单张显卡与真机工时限制 & 三种消融、三个训练种子及功效扩充的周期较长 & 复用同一预训练检查点，四台机器人区组并行；预试验后用实测链时长核算机器人小时 & 保留主比较的三个种子与样本下限，延长实验时段并削减不影响主要假设的探索性组合 \\
  四台机器人本体差异 & 同一策略在不同机器人上的成功率与不一致度分布不同 & 统一标定与零位校验；源机器人和目标机器人分组报告 & 主实验只允许阈值校准；按机器人微调仅列为主实验后的补充，不混入主结果 \\
  VR示范采集吞吐不足 & 随机停靠有效示范少于每子任务50回合 & 自由空间段可由导航与预设航点自动完成，仅遥操作精细段\cite{mandlekar2023hitltamp} & 延长采集窗口或减少探索性自然停靠回合，不降低主消融的数据量下限 \\
  \bottomrule
\end{longtable}

\section{参考文献}

\begingroup
\let\flushbottom\relax
  \bibliographystyle{hithesis}
  \bibliography{reference}
\endgroup

% Local Variables:
% TeX-master: "../report"
% TeX-engine: xetex
% End:
